%!TEX root = ../thesis.tex
%*******************************************************************************
%****************************** Third Chapter **********************************
%*******************************************************************************
\chapter{Dopamine neurons encode trial-by-trial
subjective reward value in an auction-like task}

% **************************** Define Graphics Path **************************
\ifpdf
    \graphicspath{{Chapter4/Figs/Raster/}{Chapter4/Figs/PDF/}{Chapter4/Figs/}}
\else
    \graphicspath{{Chapter4/Figs/Vector/}{Chapter4/Figs/}}
\fi

Once satisfied that the two monkeys could perform the BDM auction task and provide us with a measure of utility (the bids made), the aim of the project was to then record from the dopaminergic cells of the primate midbrain to compare this behavioural measure with a neuronal one. Dopaminergic cells have been implicated in encoding the utility offered to the monkey in the form of binary choices \cite{staufferDopamineRewardPrediction2014}. The aim of this project was to see if and how the activity of these cells reflected the utility reported by the monkey's bidding behaviour in the auction task presented in Chapter 3 (Figure \ref{fig:task_epochs}).

\clearpage
\section{Materials and Methods}
In addition to the systems previously described that generated, presented, and responded to the animal's behaviour in order to gather the behavioural data, we also had a separate system running concurrently to record and store in real time electrophysiological signals generated by neurons in the monkey's brains. Here, we detail how we accessed and recorded these signals.

\subsection{Implantation surgeries}
In order to access the central nervous system of each monkey, a cranial window was implanted in aseptic surgeries under general isoflurane anaesthetic, performed by us and assisted by a veterinary anaesthetist as well as other members of the laboratory.

For Monkey U, two implantation surgeries were performed in this regard. First, a head-fixation post (Crist instruments) was implanted into the skull using bone screws and dental cements (Super-Bond C\&B, Sun Medical; Paladur, Kulzer). This allowed for the monkey's head to be fixed in place for recording experiments. Training of the monkey in the behavioural task with head fixation was performed prior to a second surgery (see Figure \ref{fig:monkey_trials} and Section \ref{sec:bidding_rank_ordered}) to habituate the monkey to work with his head fixed in place. A second aseptic surgery was performed to stereotaxically implant a titanium recording chamber (Gray Matter Research) positioned on top of the skull (see Figure \ref{fig:xrays}) using similar bone screws and dental cement.

For Monkey V, both a similar metal chamber (Figure \ref{fig:xrays}) and bars that acted to hold the head in place were implanted in a single aseptic surgery that otherwise used similar methods. Head fixation training then took place roughly a month after this surgery. For both monkeys, a Kopf stereotaxic manipulator was used to precisely position the cranial windows.
Two weeks after implantation of the metal chamber, a craniotomy was performed on each monkey using a surgical drill to open the skull in a location estimated to be directly above the desired recording location of midbrain dopaminergic neurons. Craniotomies were monitored, maintained with tissue growth removal, and cleaned daily with saline.

\begin{figure}[htbp!] 
\centering    
\includegraphics[width=1.0\textwidth]{Chapter4/Figs/x_ray_fig}
\caption[X-ray images were used to confirm the location of the recording chamber.]{\textbf{X-ray images were used to confirm the location of the recording chamber.} X-rays for both monkeys (Monkey U left; Monkey V right) were taken post-surgery to implant the recording chambers. A cannula was placed with the tip roughly 2cm below the animal's dura mater at recorded coordinates. The relative location of this cannula to the ears and sphenoid bone of the monkey was used to judge the location of the implant and future recording coordinates.}
\label{fig:xrays}
\end{figure}

\subsection{Identifying dopaminergic neurons} \label{subsec:identification}

After access to the central nervous system was provided by the craniotomies, we located dopaminergic cells in the following 4 steps:

\begin{enumerate}
    \item A metal cannula was inserted at recorded coordinates relative to the chamber into the brain of each animal and it's position relative to landmarks was imaged using an X-ray (Siemens AG), with the aid of the Anatomy Visual Media Group in the department). Particularly useful were the locations of the sphenoid bone on the lateral edge of the skull and the metal bars placed at the location of the monkey's ears (highlighted in Figure \ref{fig:xrays}).
    \item From these landmarks, the location of the ventral posteromedial nucleus (VPM) of the thalamus was estimated. Under anaesthesia, single electrodes (see following sections) were advanced through the cortex to try and find this nucleus, which is densely packed with large and responsive neurons which encode somatosensory stimuli originating from the skin and other external tissues. These neurons exhibit stereotypical bursts when areas of the face and mouth are stimulated by touch (\cite{loeBodyRepresentationVentrobasal1977}) and so are easy to identify once reached\footnote{In practice, searching for this nucleus generally also involves searching for the adjacent ventral posterolateral nucleus of the thalamus which has similar somatotopic encoding for the arms, hands, legs, and feet of the animal and narrowing in on the VPM topographically from there.}.
    \item Once the VPM has been mapped to coordinates relative to the recording chamber, then these neurons can also be located in the awake animal. Given the large receptive fields of neurons that carry somatic information inside the mouth, such neurons could be stimulated by feeding the animal chunks of fruit. Penetrating through this nucleus a further 3mm placed the electrode in the region of dopaminergic cells of the ventral midbrain from which we could record. Dopaminergic neurons were initially identified by their location directly below the thalamus and their characteristic trace criteria: a polyphasic action potential waveform with an initial positive, followed by a sharp negative and then a longer positive component. The relatively long (>2ms) waveform duration is usually accompanied by a slow baseline firing rate at 0.5-8Hz \cite{romoDopamineNeuronsMonkey1990}, \cite{ljungbergResponsesMonkeyDopamine1992} (c.f. nearby Substantia Nigra pars reticulata neurons, which may have firing rates up to 10 times more frequent). 
    \item Finally, while searching for reward-encoding dopaminergic cells, monkeys were faced with a blank screen devoid of stimuli. We tested the responses of putative dopaminergic cells by stochastically providing the monkey with small bursts of liquid (generally 0.3ml of water) that were entirely unpredicted. See Figure \ref{fig:da_raster} for data recorded from a single dopaminergic cell in this project for the stereotypic positive reward prediction error response of dopamine neurons. To aid in the identification of putative neurons, potential traces were displayed live on an oscilloscope (TDS 2024, Tektronix) and amplified aurally.
\end{enumerate}

The locations of the recorded dopaminergic cells were verified using histological methods. In short: focal electrolytic lesions are made using 20-60s bursts of 15-20µA at coordinates and depths used to record dopaminergic cells. These lesions were then visually compared with the VTA and SNc locations by cresyl-violet Nissl staining of 40µm coronal brain sections. The monkeys were sacrificed by intravenous overdose of sodium pentobarbital and preserved using formaldehyde (4\% formalin in 0.1M phosphate buffered saline).

\subsection{Neuronal recording and data collection}
We recorded neurons using a mixture of single pre-fabricated (Alpha-Omega) and in-house custom-made electrodes. Custom-made electrodes were produced from tungsten wire etched to a point using electrolysis in a sodium nitrate/potassium hydroxide solution. This wire was then glass coated with borosilicate glass (Harvard Apparatus) using a gravity-pulling device. Under a microscope, the tip was then unsheathed from glass by melting. Finally, electrodes were first coated with gold (Neuralynx gold plating solution) and then platinum (Neuralynx 1\% chloroplatinic acid hydrate) to lower the impedance across the electrode to roughly 0.1M$\Omega$.

Before recording, the electrodes were bathed for more than 30 minutes in a 70\% ethanol solution before sterile backfeeding of the electrode through a 23 gauge stainless steel guide cannula attached to a modified micromanipulator (Monkey U: MO-97, Narishige International Ltd, Monkey V: CMS, Nan Instruments). Then, this micromanipulator was attached to an x/y stage placed on the cranial implant, which was used to precisely select coordinates relative to this implant to record from. The distribution of the coordinates recorded for both monkeys is shown in \ref{fig:recording_coords}. The electrode was advanced through the brain using manual (Monkey U) or computer-based motor-driven (Monkey V) control of the micromanipulator.

\begin{figure}[htbp!] 
\centering    
\includesvg[width=1.0\textwidth,inkscapelatex=false]{Chapter4/Figs/redone_figs/recording_heatmap.svg}
\caption[Heatmap of the recording coordinates for both monkeys.]{\textbf{Heatmap of the recording coordinates for both monkeys.} Daily recording co-ordinates relative to the centre of the recording chamber (0,0; red) were chosen based on success in previous sessions while leaving a gap of at least 1mm in both the anterior-posterior and medial-lateral dimensions from recording made in the previous session. Figure \ref{fig:recording_coords} shows a heatmap of the recording sites of all 236 neurons recorded from Monkey U and for the 259 neurons recorded from Monkey V. Size of circles represents the number of sessions on which coordinates were used for recordings and colour represents the number of cells recorded from those coordinates.}
\label{fig:recording_coords}
\end{figure}

Electrodes were sampled at 22000Hz. To amplify electrode recordings, signals were passed first through a preamplifier (custom-made, Biotronix workshop, Department of Psychology, University of Cambridge), and then successively through two further steps of amplification before bandpass filtering at 100-5000Hz and 100-3000Hz successively. Between amplification steps, line noise was reduced using a Humbug (AM Systems).

\subsection{Data analysis}

Both the amplified raw analog signal from the electrode and digitised pulses were stored via a National Instruments Data Acquisition Card (PCI-6229) and custom-built internal lab software (‘Getty') running on MATLAB 2014 (MATLAB, 2017) calling C++.

Though we recorded single units, identified visually and aurally (see previous section), some offline spike sorting was done using Spike2 software (Cambridge Electronic Design Ltd) to extract individual waveforms and remove noise and artefacts when necessary.
Cell traces are smoothed for plotting via convolution with the following function:

\begin{equation}
    (1 - \mathrm{e}^{-t}) \cdot \mathrm{e}^{-\frac{t}{T}}
    \label{eq:smoothing}
\end{equation}

Where $t$ is the time from the moving window function and $T$ is set to 200 for our purposes. Such a kernel has previously been used \cite{eshelDopamineNeuronsShare2016} in similar studies. 

To aggregate cells, the smoothed firing rate is normalised via z-scoring per cell by subtracting the mean ($\mu$) firing rate for that cell in the time immediately before a point of interest (e.g., the reveal of a trial outcome) and dividing by the standard deviation ($\sigma$):

\begin{equation}
    z = \frac{x - \mu}{\sigma}
    \label{eq:z-smooth}
\end{equation}

This standardization ensures that the observed population responses and regression slopes reflect the relative sensitivity of the dopaminergic population, rather than being skewed by neurons with high absolute firing rates or different dynamic ranges.

Data transformations were carried out in MATLAB 2020b and R 4.0.2 (‘Taking Off Again'; \cite{rcore}). Statistical tests were also done in R 4.0.2 on the same computer used for behavioural analysis. These tests were carried out using the functions \texttt{wilcox.test} and \texttt{kruskal.test} of the \texttt{stats} package which comes prebundled with modern R installations, and \texttt{dunn.test} from the package of the same name. All Wilcoxon signed-rank tests are one-sided unless otherwise specified. This is due to our prior beliefs that neuronal responses would scale positively with reward value, consistent with the known reward-coding properties of dopaminergic neurons.

\clearpage
\section[Responses and Properties of Recorded Putative Dopaminergic Neurons]{Responses and Properties of Recorded Putative\\Dopaminergic Neurons}

\subsection{Canonical responses of putative dopaminergic neurons}

Over the course of the project, we recorded responses from 268 putative dopamine neurons; 123 from Monkey U and 145 from Monkey V. The average cell trace and interspike interval is shown for 2 cells from each animal in Figure \ref{fig:isi_waveforms}, showing the characteristic long waveform (left) and slow firing rate (right) expected from dopaminergic neurons. See the text of the figure and Section \ref{subsec:identification} for precise details.

In addition to these stereotypical characteristics, the 268 putative dopaminergic neurons described showed a significant response correlated to at least one epoch of the auction task (that is, the display of the fractal, the bid made by the monkeys, or the reveal of the outcome of the trial; Figure \ref{fig:task_epochs}).

\begin{figure}[p]
\centering    
\includesvg[width=1.0\textwidth,inkscapelatex=false]{Chapter4/Figs/redone_figs/cell_waveforms_fig}
\caption[Spike-sorted putative dopamine neurons show standard characteristics.]{\textbf{Spike-sorted putative dopamine neurons show standard characteristics.} The average waveform of two isolated neurons (Spike2 software) are shown for both Monkey U (blue) and Monkey V (orange) on the left. Both waveforms show the standard polyphasic and long (greater than 2ms) characteristics of midbrain dopamine neurons. Grey ribbon shows standard deviation of potential. The inter-spike interval histogram is shown for both pairs (right) of neurons. Dopamine neurons fire relatively sparsely which is mirrored in the behaviour of these cells which show a Poisson-like interspike-interval histogram with a mean of 200-350ms between spikes. Inversely, the baseline firing rate for the four neurons pictured over the whole recording period were 3.960Hz, 3.067Hz, 3.654Hz, and 4.710Hz respectively.}
\label{fig:isi_waveforms}
\end{figure}

\subsection{Responses of putative dopaminergic neurons to trial onset}

Given the high reward rate (and therefore potentially limited trial numbers) of the auction task, we had to be reasonably confident that a single neuron picked up by an electrode was a reward-encoding dopaminergic neuron. To test this, we examined the stereotypical reward responses (\cite{mirenowiczImportanceUnpredictabilityReward1994}) to unpredicted, valuable stimuli characteristic of midbrain dopaminergic neurons before allowing the monkey to complete any trials. We therefore would stochastically release a small, 0.3ml water reward from the setup once a neuron of sufficient confidence and recording quality was captured. The response of the liquid-deprived monkey's putative dopaminergic cell to the slight quenching of its thirst was then used as an arbitrary, on-line criteria for recording the neuron within the task. As previously mentioned, despite trying to record extracellularly from single cells, the electrode would pick up potentials from other nearby neurons. These were then later sorted offline.

The dopaminergic cell population response to unpredicted liquid reward was unsurprisingly significant for both monkeys (Wilcoxon signed-rank, Monkey U: p < 0.0001, Monkey V: p < 0.0001). However, testing of cells with free reward was limited to avoid satiation of the monkey and this for many cells where only a few such unpredicted rewards were given and responses measured. When restricted to only cells that were tested with free rewards at least 5 times, 79.6\% (78/98) of Monkey U cells and 87\% (20/23) of Monkey V cells were significant in a Wilcoxon signed-rank at p < 0.05\footnote{C.f. Figure \ref{fig:free_reward} where numbers denote the same test carried out on \textit{all} putative dopaminergic cells, yielding 80/123 and 33/145 cells for Monkey U and Monkey V respectively.}.

The responses of individual neurons after sorting are shown in Figure \ref{fig:free_reward_psth}, where Figure \ref{fig:free_reward_psth}A) shows the PSTH of putative dopaminergic neurons from Monkeys U and V, respectively, responding to an unpredicted reward. This offline sorted also yielded putative dopaminergic neurons which did not respond to task epochs (Monkey U: 16/139 neurons; Monkey V: 19/164 neurons). Individual examples of PSTH for potentially dopaminergic neurons which did not respond to an unpredicted reward are shown for Monkey U and Monkey V in Figure \ref{fig:free_reward_psth}B). Finally, the sorting also produced cells that did not meet our classification criteria for dopaminergic neurons and did not respond to a free reward stimulus (Monkey U, total n= 97; Monkey V, total n = 95; Figure \ref{fig:free_reward_psth}C).

For an additional plot of the unpredicted reward response to a single such stimulus by a neuron within this thesis, see also Figure \ref{fig:da_raster}.

\begin{figure}[p]
\centering    
\includesvg[width=1.0\textwidth,inkscapelatex=false]{Chapter4/Figs/redone_figs/free_reward_psth_fig_redone.svg}
\caption[Individual neuronal responses to unpredicted reward.]{\textbf{Individual neuronal responses to unpredicted reward.} Putative dopaminergic neurons were identified via their baseline firing rate and polyphasic waveform. After identifying such neurons, an unpredicted reward of 0.3ml water was delivered to the monkeys to test the canonical dopaminergic reward response. Responses from Monkey U are shown in blue on the left and Monkey V in orange on the right. Spike counts within each panel are normalised to the maximum bin count for that monkey, giving a common 0--1 y-axis scale; absolute spike counts therefore differ between panels. Grey boxes indicate time windows from which impulses were counted for statistical tests: 1.5--0.5s prior to the rewarding stimulus (baseline) and in the 500ms immediately after (red dashed line). Similar denoting of all analysis windows is maintained for all figures in this chapter. A) shows the PSTH of putative dopaminergic neurons that showed a significant (Wilcoxon signed-rank test, p $<$ 0.05) response to unpredicted reward (113/268 recorded neurons). B) shows examples of putative dopaminergic neurons that did not show a significant response to unpredicted reward for Monkey U and Monkey V. Finally, post-recording sorting yielded putative non-dopaminergic neurons for both monkeys (n = 113 Monkey U, n = 114 Monkey V); a lack of response to unpredicted reward for one such cell per monkey is shown in C).}
\label{fig:free_reward_psth}
\end{figure}

The responses of all 268 putative dopaminergic neurons to these bursts of liquid are shown in Figure \ref{fig:free_reward}, with both the average standardised firing rate of these neurons per monkey (Figure \ref{fig:free_reward}A) and a raster plot of each suprathreshold action potential spike of neurons with at least 5 free reward stimuli shown below (Figure \ref{fig:free_reward}B) with each row showing every response to an unpredicted reward encountered by an individual neuron superimposed on across multiple free water reward stimuli. The raster in Figure \ref{fig:free_reward}B is restricted to cells tested with more than 5 unpredicted rewards that also showed a significant response (Wilcoxon signed-rank, p < 0.2); this lenient threshold was used as a practical inclusion criterion rather than a formal statistical claim. The remaining putative dopaminergic neurons contribute to the population average in A) but do not appear in the raster because free reward testing was deliberately limited to avoid satiation, leaving the Wilcoxon test underpowered for cells given only a small number of stimuli. Among cells tested sufficiently (n > 5), approximately 80\% showed a significant response for both monkeys, consistent with the majority of recorded cells being reward-responsive.

\begin{figure}[p]
\centering
\includegraphics[width=0.98\textwidth]{Chapter4/Figs/redone_figs/free_reward_fig.png}
\caption[Putative dopaminergic neurons show reward response behaviour to unpredicted reward.]{\textbf{Putative dopaminergic neurons show reward response behaviour to unpredicted reward.} The z-normalised firing rate of all neurons categorised as reward-encoding putative dopaminergic neurons for Monkey U (A) left; n = 123 neurons) and Monkey V (A) right; n = 145 neurons) are shown. Activity in the 0-500ms following the delivery of 0.3ml unpredicted water reward (dashed line) is highly significant for both populations on a Wilcoxon signed-rank test (Monkey U, p < 0.0001; Monkey V, p < 0.0001). The responses of neurons showing significant (Wilcoxon signed-rank test) responses to an unpredicted reward shown below on a raster plot of individual neuronal impulses for Monkey U (n = 80 neurons, B) left) and Monkey V (n = 33 neurons, B) right).}
\label{fig:free_reward}
\end{figure}

During recording sessions, we noticed that the putative dopaminergic cells we recorded showed robust activation by a trial onset signal shown at the beginning of each trial to indicate that the monkey must keep the joystick centred (see Chapter 3). There is a random 2.5 to 3.5s interval between the final epoch of a trial (the budget water payout) and the signalling of the start of the next trial (see Figure \ref{fig:task_epochs}). The responses to the signal of onset for the cell population included in our analysis (Monkey U: n = 123 neurons, Monkey V: n = 145 neurons) are shown in Figure \ref{fig:fixation_cross}.

Normalised population firing rate (Figure \ref{fig:fixation_cross}A) and individual cell raster plots (Figure \ref{fig:fixation_cross}B) are shown for all of these cells per monkey. For Monkey U, 82 out of 123 neurons showed significant responses to the fixation cross signal at the beginning of each trial (Wilcoxon signed-rank test, p < 0.05). For Monkey V, 62 out of 145 neurons showed significant responses on the same test. The grey box prior to the onset of signal (red dotted line) encapsulates the 500ms prior to this unpredicted trial start signal and is used as a baseline for both this and future figures and statistical tests of population activity.

\begin{figure}[p]
\centering    
\includegraphics[width=0.95\textwidth]{Chapter4/Figs/redone_figs/fixation_cross_fig.png}
\caption[Putative dopaminergic neurons are sensitive to the onset of trials.]{\textbf{Putative dopaminergic neurons are sensitive to the onset of trials.} The z-normalised firing rate (Monkey U left; Monkey V right) of putative dopaminergic neurons (n = 123 Monkey U, n = 145 Monkey V) show an increased activity in response to the display of a fixation cross indicating the start of each trial (dashed line). During this time period, the monkey must hold the joystick steady. The increase in neuronal activity in the 100-600ms following the onset stimulus was tested using a signed-rank test (p < 0.05) of spike rate compared to the 500ms immediately preceding the stimulus (dashed red line). 82/123 neurons showed significant increases in activity for Monkey U and 62/145 neurons showed significant increases in activity for Monkey V. Raster plots of each of these neuron's responses at this epoch are shown in B) (Monkey U left; Monkey V right).}
\label{fig:fixation_cross}
\end{figure}

This response to the onset of the trial is not surprising. As discussed previously, individual auction trials are highly rewarding, and monkeys can expect to win >1ml of liquid on an average trial. We hypothesise therefore that for an attentive monkey, the signal that they will be allowed to complete a trial with only reasonable cognitive and motor effort required is a conditioned stimulus for a later predicted reward (either juice or liquid)\footnote{See Section \ref{subsec:hedonimetry} (‘Hedonimetry and the Brain') for more on dopaminergic responses during conditioning.}.

\clearpage
\section{Dopaminergic Responses Signal Subjective Reward Value}

Although BDM trials are on average highly rewarding, there are two pieces of information within each trial that indicate just how much liquid a single trial will deliver to a monkey. The first of these occurs 1.2 seconds into each trial, where a fractal is used to indicate the volume of juice reward the monkey is likely to win should they outbid the computer (see Figure \ref{fig:task_epochs}).

\subsection{Responses of putative dopaminergic neurons to display of reward signalling fractal}

Figure \ref{fig:fractal_display} shows the responses of the recorded putative dopaminergic neurons to the display of the fractal indicating either a low, medium, or high value reward on offer. Initially, the population activity for all three offer-signalling fractals is increased, though this is only sustained for trials in which the fractal indicating the offer of the highest value fractal is shown to the monkey. This ‘two component' dopaminergic prediction error signalling has previously been described in the literature (\cite{schultzDopamineRewardPredictionerror2016} for a review; \cite{nomotoTemporallyExtendedDopamine2010}, \cite{staufferDopamineRewardPrediction2014}). The first, non-selective, component signalling the sensory detection and attention toward a salient reward-associated stimuli compared to a later second component associated with valuation processing. Only this second component of the signal shows bidirectional responses, i.e. a highly rewarding stimuli will elicit greater neuronal responses whereas a disappointing auctioned good may lead to a depression in activity.

\begin{figure}[htbp!] 
\centering    
\includesvg[width=1.0\textwidth,inkscapelatex=false]{Chapter4/Figs/redone_figs/fractal_display_fig}
\caption[Putative dopaminergic neurons show ranked-ordered activity to the offered good.]{\textbf{Putative dopaminergic neurons show ranked-ordered activity to the offered good.} Averaged normalised responses for all putative dopaminergic neurons recorded from sessions with three reward levels (n = 102 Monkey U; n = 121 Monkey V) show an initial increase in firing rate within the first 100ms of the offered fractal good being revealed before showing rank ordered activity in the second response component (grey box). Over the population of putative dopaminergic cells, a Kruskal-Wallis test was significant for the effect of fractal offer (Monkey U, p < 0.0005; Monkey V, p < 0.0001).  The difference between pre-trial and second component activity is also tested per individual cell using a linear regression of neuronal response to the discrete reward value offer. N = 31/102 (Monkey U) and n = 41/121 (Monkey V) putative dopamine neurons were significant at p < 0.05 for rank ordered responses under this test.}
\label{fig:fractal_display}
\end{figure}

A Kruskal-Wallis test on the proportion of spikes in a defined window period after the fractal reveal, grouped by the reward level of these stimuli, is strongly significant, with maximally increased putative dopaminergic impulses for the most potentially rewarding offer (Monkey U: p < 0.0001, Monkey V: p < 0.0001; Figure \ref{fig:fractal_display}). This window was defined per monkey. We took the mean baseline population response of all putative dopaminergic neurons per money as the activity in the 500ms before the trial-onset fixation cross stimuli. The start of the second component analysis window began when the population activity (split into 20ms bins) of neurons responding to the display of the offered fractal good was no longer significantly different from the baseline (Bonferroni-Holm corrected Wilcoxon signed-rank test, p < 0.05). The end of the window was defined when the 20ms bins of normalised neuronal activity for the highest reward fractal for the monkey returned to the baseline using the same Bonferroni-Holm corrected Wilcoxon signed-rank test. A sketch of this process is described in Figure \ref{fig:analysis_windows}.

\begin{figure}[htbp!] 
\centering    
\includegraphics[width=0.9\textwidth]{Chapter4/Figs/analysis_window.png}
\caption[Analysis windows for dopaminergic response to offered good is defined per monkey.]{\textbf{Analysis windows for dopaminergic response to offered good is defined per monkey.} Due to the two-component nature of dopaminergic signalling, we defined the window over which statistics were carried out on the fractal offer epoch per monkey. Firstly, the average activity in the 500ms leading up to the stochastic display of the fixation cross (indicating the onset of a trial) was taken as a baseline of activity (first green rectangle). 1.2s after this epoch, a fractal is displayed indicating the offered juice reward magnitude for the trial. The initial dopaminergic response to this (red lines) is positive, indicating attention to the stimulus. This signal is split into 20ms bins and each is compared to the baseline signal using a Bonferroni-Holm signed-rank test at p < 0.05 significance. Using the lowest value reward stimulus, we expect a negative second component RPE, therefore, we take the point at which the dopaminergic signal is no longer significantly different from baseline as the onset of the analysis window (second green rectangle). The end of the analysis window is then defined at which the dopaminergic RPE signal again returns to baseline following an (expected) reduction is cellular activity. For Monkey U and Monkey V this window was defined as 180-340ms post stimulus and 180-360ms post stimulus respectively. C.f. Figure \ref{fig:task_epochs}.}
\label{fig:analysis_windows}
\end{figure}

The magnitudes of potential reward offer increase linearly (0.2, 0.45, 0.7ml for Monkey U; 0.3, 1.0, 1.7ml for Monkey V), and so we might expect a linear change in normalised putative dopaminergic impulses in response: symmetric increases and decreases in firing rate for the corresponding high and low juice offers, with little change in firing rate for the fractal which represents the average pre-trial offer. This is not what we see with small decreases in the number of measured action potentials for both low and medium value reward offers (after an initial first-component positive signal for both), and much larger increases in firing rate when the fractal indicating a higher valued offered good is shown (Figure \ref{fig:fractal_display}). The valuation component of the dopaminergic signal seems to treat the offer of a lesser-valued good as a loss, which has an outsized ‘hurt' compared to the potential gain of the high-valued fractal \cite{kahnemanProspectTheoryAnalysis1979}.

As the magnitude of the reward offer and the subsequent bid made by the monkey would be expected to be somewhat correlated, we also recorded cells from a number of sessions in which only the middle reward offer was presented to each monkey on each trial. Therefore, any variability in bid made is only reliant on the subjective value the monkey has for the taste/nutrition of the fruit juice compared to the water budget. Figure \ref{fig:one_fractal_display} shows the normalised population response of 45 putative dopaminergic neurons\footnote{Neurons were categorised as putative dopaminergic neurons using the same procedure described in Figures \ref{fig:isi_waveforms} and \ref{fig:free_reward_psth}} recorded from sessions where only one offer magnitude was shown to the monkey (Monkey U: n = 21 neurons; Monkey V: n = 24 neurons). 

Monkeys U and V showed opposite averaged population activity to the presentation of the offered middle reward fractal, with Monkey U showing a decrease in second-component response and Monkey V showing an increase. Compare this with Figure \ref{fig:fractal_display} where a similar response is observed per monkey to the middle reward offer fractal in sessions where all three reward magnitudes are shown. For the analysis of sessions comprising single fractal offers, the same second component analysis windows as described in Figure \ref{fig:analysis_windows} are used. In a two-sided Wilcoxon test, 3 putative dopaminergic neurons for Monkey U and 7 putative dopaminergic neurons recorded from Monkey V showed a significant and consistent valence of responses to the display of the offered good.

\begin{figure}[htbp!] 
\centering    
\includesvg[width=1.0\textwidth,inkscapelatex=false]{Chapter4/Figs/redone_figs/one_fractal_display_fig}
\caption[Normalised dopaminergic responses when only one fractal is shown per trial block.]{\textbf{Normalised dopaminergic responses when only one fractal is shown per trial block.} When only the middle reward magnitude was shown in a trial blocks, average population activity showed opposite valence between monkeys U (decrease in firing activity) and V (increase). Per monkey, 3/21 (Monkey U) and 7/24 (Monkey V) cells recorded in trial blocks where only one reward fractal level was shown had significant responses (two-sided Wilcoxon, p < 0.05).}
\label{fig:one_fractal_display}
\end{figure}

The average responses of individual cells that showed a significant response (see Figure \ref{fig:fractal_display}) to the display of the fractal offer are shown in Figure \ref{fig:fractal_display_raster}. The averaged responses of 31 cells recorded from Monkey U (left) and 41 cells recorded from Monkey V (right) are shown per trial, separated by the level of the reward offered. The top panel shows the averaged responses on the trials in which the highest reward magnitude was offered, the middle where the medium reward magnitude was offered, and the bottom where the lowest reward magnitude was offered.

\begin{figure}[htbp!] 
\centering    
\includegraphics[width=1.0\textwidth]{Chapter4/Figs/Vector/fd_response_raster_fig}
\caption[Responses of individual rank-order responsive neurons to each fractal offer.]{\textbf{Responses of individual rank-order responsive neurons to each fractal offer.} The normalised average responses of each individual neuron (single row) to either the high (top), medium, or low (bottom) fractal are shown for neurons which showed significant (linear regression, p < 0.05) rank ordering of second component activity to the offered good shown.}
\label{fig:fractal_display_raster}
\end{figure}

\subsection{Dopamine neurons encode trial-by-trial
subjective reward value in an auction-like task}

The central idea in this thesis is not simply that dopamine neurons encode the magnitude of the physical reward offered to them (a well-characterized finding, see \cite{toblerAdaptiveCodingReward2005}; \cite{fiorilloDiscreteCodingReward2003}), but that neurons encode the utility of the good offered to the animal. After the display of the fractal offer, the monkey is given four seconds in which to make a bid indicating their utility for the juice reward on offer. In the previous chapter, we showed that monkeys behave as if they are doing this when placing bids (see Figure \ref{fig:optimal_bidding}); now we question whether the activity of putative dopaminergic neurons of their midbrain mirrors this.

When responses to the display of the offered reward magnitude are grouped by the fifth of the bidspace within which the monkey ends up bidding, it shows a clear separation of the second component dopaminergic activity (Figure \ref{fig:fractal_display_quintiles}). A Kruskal-Wallis test for the difference in median population activity was significant for both monkeys (Monkey U p < 0.005; Monkey V p < 0.0001) indicating that these neurons are not just encoding the magnitude of the juice being offered to the monkey, but the animals own, internal, subjective value for the good.

At the level of individual neurons, we studied the relationship between the second component firing rate and the subsequent monkey bid using a simple linear regression described in Equation \ref{eq:firing_rate_mb_regression}\footnote{N.b. for this and all regressions in this chapter, trials in which the monkey bid 0ml of water are included in these regressions despite these implying that either the monkey feels there is no value on offer, or that they are not engaging properly with the task.}.

\begin{align}
\text{Firing Rate} &\sim \beta_0 + \beta_1(\text{Monkey Bid}) + \varepsilon \label{eq:firing_rate_mb_regression}
\end{align}

Under this regression, we identified 73 cells (Monkey U 32/123; Monkey V 41/145) with a significant correlation (p < 0.05) of activity upon offering versus the eventual bid for a juice reward-indicating fractal.

\begin{figure}[htbp!] 
\centering    
\includesvg[width=1.0\textwidth,inkscapelatex=false]{Chapter4/Figs/redone_figs/fractal_display_bids_fig.svg}
\caption[Putative dopaminergic neurons activity at the fractal offer is correlated to monkey's final bid.]{\textbf{Putative dopaminergic neurons activity at the fractal offer is correlated to monkey's final bid.} Averaged normalised responses for all putative dopaminergic neurons (n = 123 Monkey U; n = 145 Monkey V) coloured by the quintile of the total bidding space (5 is highest, 1 is lowest) of the monkey's forthcoming bid. For both monkeys, trials on which a higher final bid will be made show increases in putative dopaminergic cell activity, while trials where the monkey makes a lower final bid show decreases in activity post juice offer reveal. The population of putative dopaminergic cells showed significant second-component activity correlated with bid quintile (Kruskal-Wallis test; Monkey U, p < 0.005, Monkey V, p < 0.0001). A simple linear regression of continuous monkey bid against putative dopamine neuronal response yielded 32/123 (Monkey U) and 41/145 (Monkey V) cells that showed significant (p < 0.05) correlation.}
\label{fig:fractal_display_quintiles}
\end{figure}

The normalised responses of these cells are averaged over all trials regardless of the fractal offered and plotted in Figure \ref{fig:monkey_bid_fractal_display_raster}

\begin{figure}[htbp!] 
\centering   
\includegraphics[width=1.0\textwidth]{Chapter4/Figs/Vector/monkey_bid_fd_raster_fig}
\caption[Responses of individual neurons sorted by forthcoming monkey bid.]{\textbf{Responses of individual neurons sorted by forthcoming monkey bid.} The normalised average responses of each individual neuron (single row) to the display of the fractal offer are shown, sorted by the fifth of the bidspace the forthcoming monkey bid is placed within. Only putative dopaminergic neurons which showed significant (linear regression, p < 0.05) correlation in activity to the monkey bid (n = 31/123 Monkey U; n = 41/145 Monkey V) are shown.}
\label{fig:monkey_bid_fractal_display_raster}
\end{figure}

To study the correlation between population activity and forthcoming bid, we take the putative dopaminergic cells (Monkey U n = 123; Monkey V n = 145) and bin bids per trial by the $25^{\text{ths}}$ of the total bidding space (i.e., into bins of 0.048ml) they fall into, irrespective of the fractal offered. In Figure \ref{fig:fractal_display_regression}A) the average firing rate ($\pm$ standard error of the mean) of putative dopaminergic neurons in the analysis window after the display of the fractal is regressed against these bins (Monkey U blue, left; Monkey V orange, right). In Figure \ref{fig:fractal_display_regression}B) the firing of neurons in the analysis window is normalised by the average response of each cell prior to the onset of the trials. This is regressed against the same binned data as in A). The formulae for these linear regressions are described in Equations \ref{eq:firing_rate_regression} and \ref{eq:normalized_firing_rate_regression}.

\begin{align}
\text{Firing Rate} &\sim \beta_0 + \beta_1(\text{Monkey Bid Bin}) + \varepsilon \label{eq:firing_rate_regression} \\
\frac{\text{Firing Rate}}{\text{Baseline Activity}} &\sim \beta_0 + \beta_1(\text{Monkey Bid Bin}) + \varepsilon \label{eq:normalized_firing_rate_regression}
\end{align}

Both preparations of the data in Figure \ref{fig:fractal_display_regression} show a significant correlation of the activity of putative dopaminergic neurons on the final offer made in these linear regressions. For details on the results, see Table \ref{tab:fd_regression_results}. It should be noted that given the reaction times of monkeys (Table \ref{tab:monkey_movement}), these neuronal responses correspond to a period of time before the monkey has, on average, begun to place their bid.

The regressions above pool trials as though they are independent observations, when in fact roughly 100 trials are recorded from each neuron, and neurons are recorded from only two independent monkeys. Data with this kind of hierarchical structure risk inflating the apparent significance of a fixed-effect estimate if this repeated-measures structure is not accounted for. To address this, we additionally fitted a linear mixed-effects model with a random intercept per cell, using \texttt{lme4}/\texttt{lmerTest} (\cite{bates2015fitting}; \cite{kuznetsova2017lmertest}):

\begin{align}
\frac{\text{Firing Rate}}{\text{Baseline Activity}} = \beta_0 + \beta_1(\text{Monkey Bid}) + u_{\text{cell}} + \varepsilon, \qquad u_{\text{cell}} \sim \mathcal{N}(0, \sigma^2_{\text{cell}}) \label{eq:lme_general}
\end{align}

where $\beta_1$ is the estimated coefficient of the fixed-effect predictor of interest, and $u_{\text{cell}}$ is a random intercept per cell absorbing cell-to-cell baseline variability, accounting for the repeated-trials-within-cell structure of the data, rather than being a fixed effect itself. Each cell's baseline offset is assumed to be normally distributed around zero, with a single shared variance term describing how much cells typically differ from one another.\footnote{All regressions in this chapter, including this mixed-effects model, are fitted separately per monkey rather than pooling across animals.}

This mixed-effects model confirmed a significant positive relationship between bid and normalised firing rate for both monkeys, fitted directly on trial-level data without binning (Monkey U: $\beta = 0.781 \pm 0.070$, $t = 11.2$, $p = 9.70 \times 10^{-29}$; Monkey V: $\beta = 0.644 \pm 0.053$, $t = 12.3$, $p = 1.68 \times 10^{-34}$). This confirms that the central bid--dopamine relationship reported above holds once the hierarchical, repeated-measures structure of the data (multiple trials per cell) is properly accounted for, rather than being an artefact of treating non-independent trials as though they were independent observations.

\begin{table}[htbp]
\centering
\begin{tabular}{lccccc}
\hline
\textbf{Monkey} & \textbf{Model} & \textbf{Intercept} & \textbf{Slope} & \textbf{Adj $R^2$} & \textbf{$p$-value} \\
\hline
Monkey U & Absolute & $1.674 \pm 0.090$ & $1.760 \pm 0.151$ & 0.850 & $3.70 \times 10^{-11}$ \\
Monkey U & Normalised & $-0.369 \pm 0.037$ & $0.671 \pm 0.062$ & 0.829 & $1.63 \times 10^{-10}$ \\
Monkey U & Mixed-Effects & $-0.426 \pm 0.054$ & $0.781 \pm 0.070$ & --- & $9.70 \times 10^{-29}$ \\
\hline
Monkey V & Absolute & $1.467 \pm 0.120$ & $2.677 \pm 0.201$ & 0.880 & $2.66 \times 10^{-12}$ \\
Monkey V & Normalised & $-0.350 \pm 0.039$ & $0.622 \pm 0.066$ & 0.784 & $2.50 \times 10^{-9}$ \\
Monkey V & Mixed-Effects & $-0.377 \pm 0.050$ & $0.644 \pm 0.053$ & --- & $1.68 \times 10^{-34}$ \\
\hline
\end{tabular}
\caption[Putative dopaminergic neuron activity was significantly correlated with forthcoming monkey bid.]{\textbf{Putative dopaminergic neuron activity was significantly correlated with forthcoming monkey bid.} Results from linear regression described in Equations \ref{eq:firing_rate_regression} and \ref{eq:normalized_firing_rate_regression}. Binned forthcoming bid on the BDM task was significantly correlated with both the absolute and normalised spiking activity in putative dopaminergic neurons of monkey midbrain. N = 123 neurons (Monkey U) and n = 145 neurons (Monkey V). Mixed-Effects rows show the linear mixed-effects model described in Equation \ref{eq:lme_general}, fitted on trial-level (unbinned) data with a random intercept per cell; Adj $R^2$ is not directly applicable to this model type and is left blank.}
\label{tab:fd_regression_results}
\end{table}

\begin{figure}[htbp!]
\centering
\includesvg[width=0.9\textwidth,inkscapelatex=false]{Chapter4/Figs/redone_figs/fractal_display_lme_regression_fig}
\caption[Trial-level mixed-effects model predictions of dopaminergic activity against monkey bid.]{\textbf{Trial-level mixed-effects model predictions of dopaminergic activity against monkey bid.} Predicted normalised firing rate from the linear mixed-effects model described in Equation \ref{eq:lme_general}. Each grey line shows one cell's predicted response, combining the fixed-effect slope with that cell's own random intercept; the coloured line shows the population fixed-effect (Monkey U, blue, left; Monkey V, orange, right). Unlike Figure \ref{fig:fractal_display_regression}, bid is entered continuously at the trial level throughout, with no binning at any stage.}
\label{fig:fractal_display_lme_regression}
\end{figure}

\begin{figure}[htbp!] 
\centering    
\includesvg[width=0.945\textwidth,inkscapelatex=false]{Chapter4/Figs/redone_figs/fractal_display_regression_fig}
\caption[Putative dopaminergic neurons activity at the fractal offer predicts monkey's forthcoming bid.]{\textbf{Putative dopaminergic neurons activity at the fractal offer predicts monkey's forthcoming bid.} Activity following the display of the fractal reward is grouped by the $25^{\text{th}}$ of the bidspace the forthcoming monkey bid is placed within for putative dopaminergic neurons. Average firing rate per cell is shown in A) per monkey (Monkey U, n = 123 neurons, blue, left; Monkey V, n = 145 neurons, orange, right) with error bars for standard error of the mean. There is significant correlation between forthcoming monkey bid and the second component dopaminergic response (simple linear regression; Monkey U, p < 0.0001, adjusted  $R^2$ = 0.85; Monkey V, p < 0.001, adjusted $R^2$ = 0.88). Shaded area shows 95\% confidence interval of regression. B) when activity is normalised by the cell's baseline activity prior to the start of the trial, significant correlation remains (simple linear regression; Monkey U, blue, p < 0.0001, adjusted $R^2$ = 0.83; Monkey V, orange, p < 0.001, adjusted $R^2$ = 0.78).}
\label{fig:fractal_display_regression}
\end{figure}

The activity of dopamine neurons has previously been shown to be modulated by motivation rather than movement (\cite{ljungbergResponsesMonkeyDopamine1992}; \cite{satohCorrelatedCodingMotivation2003}), nevertheless we also investigated the effect of cellular firing on the gross and net (see Figure \ref{fig:monkey_joystick1}) bidding movement and also the time spent bidding by the monkey. To control for the eventual monkey bid which is highly correlated with these variables, we regress these movement parameters against the residuals of the regression in Equation \ref{eq:firing_rate_mb_regression}. The regression therefore takes the form \ref{eq:residual_regression} where the 'Movement Parameter Bin' is each of the gross bid marker movement, net bid marker movement, or total bidding time binned into $25^{\text{ths}}$ of the parameter values per monkey.

\begin{align}
\text{Residual}(\text{Firing Rate}) &\sim \beta_0 + \beta_1(\text{Movement Parameter Bin}) + \varepsilon \label{eq:residual_regression}
\end{align}

None of these parameters are significantly correlated with the residuals of the regression of monkey bid against dopaminergic firing rate (simple linear regression, p < 0.05). The complete results for these regressions are shown in Table \ref{tab:residual_regression_results} and graphically in Figure \ref{fig:residual_movement_regressions}.

\begin{table}[htbp]
\centering
\begin{tabular}{lccccc}
\hline
\textbf{Monkey} & \textbf{Movement Parameter} & \textbf{Intercept} & \textbf{Slope} & \textbf{Adj $R^2$} & \textbf{$p$-value} \\
\hline
Monkey U & Gross Movement & $0.007$ & $-0.051 \pm 0.095$ & $-0.030$ & $0.592$ \\
Monkey U & Net Movement & $-0.027$ & $0.086 \pm 0.088$ & $-0.002$ & $0.341$ \\
Monkey U & Movement Time & $0.050$ & $-0.056 \pm 0.082$ & $-0.023$ & $0.506$ \\
\hline
Monkey V & Gross Movement & $0.071$ & $0.005 \pm 0.168$ & $-0.043$ & $0.977$ \\
Monkey V & Net Movement & $0.128$ & $-0.133 \pm 0.086$ & $0.054$ & $0.138$ \\
Monkey V & Movement Time & $0.065$ & $-0.037 \pm 0.073$ & $-0.032$ & $0.618$ \\
\hline
\end{tabular}
\caption[Movement parameters were not significantly correlated with the residuals of firing rate against monkey bid.]{\textbf{Movement parameters were not significantly correlated with the residuals of firing rate against monkey bid.} Results from linear regression described in Equations \ref{eq:residual_regression}. Binned joystick movement parameters on the BDM task were not significantly correlated with the residuals of the regression described in Equation \ref{eq:firing_rate_mb_regression}.}
\label{tab:residual_regression_results}
\end{table}

\begin{figure}[htbp!] 
\centering    
\includesvg[width=1.0\textwidth,inkscapelatex=false]{Chapter4/Figs/redone_figs/movement_residual}
\caption[Binned movement parameters show little relationship with the residuals of regressions of firing rate against monkey bid.]{\textbf{Binned movement parameters show little relationship with the residuals of regressions of firing rate against monkey bid.} Movement parameters for 'gross movement' (that is, the total movement of the bid marker over the course of a trial), 'net movement' (the difference between the starting and final bid) and 'movement time' (the time from when the monkey starts to move the bid marker to when they finally place it) are plotted against the residual of the regression in Equation \ref{eq:firing_rate_mb_regression} for all putative dopaminergic neurons (Monkey U, n = 123; Monkey V, n = 145). None of the binned movement parameters show significant relationship to the residual of this regression (simple linear regression, p < 0.05).}
\label{fig:residual_movement_regressions}
\end{figure}

\subsubsection{Dopaminergic Responses in Single-Fractal Sessions}
\label{subsec:dopamine_single_fractal}

When applying the same methods to data from sessions in which only the medium level fractal offer was shown to the monkeys, the picture is much less clear. We binned neuronal data by terciles of the bidding range instead of quintiles as in Figure \ref{fig:fractal_display_quintiles} due to the sparsity of the data from these sessions (Monkey U 2228 trials from 22 sessions; Monkey V 2724 trials from 22 sessions; See Figure \ref{fig:single_fractal_bids}), otherwise the methods are the same. The same analysis windows are used for the single fractal sessions as for the sessions where a range of goods are offered. Monkey V showed a significant ranking order of population activity within this second component window (Kruskal-Wallis test, p < 0.05) though Monkey U did not (p = 0.45).

\begin{figure}[p] 
\centering    
\includesvg[width=1.0\textwidth,inkscapelatex=false]{Chapter4/Figs/redone_figs/one_fractal_display_bids_fig}
\caption[Monkey bid correlation to dopaminergic response at fractal display bid in sessions where only one reward magnitude is offered.]{\textbf{Monkey bid correlation to dopaminergic response at fractal display bid in sessions where only one reward magnitude is offered.} When only one fractal is offered to the monkeys in a session, the normalised responses for putative dopaminergic neurons (Monkey U n = 21 cells, Monkey V n = 24 cells) are shown split by thirds of the total bidding space (3 is the top third, 2 the middle, and 1 the lowest third) where a bid is later placed. The population of putative dopaminergic cells showed significant second-component activity correlated with bid tercile for Monkey V (p < 0.05, Kruskal-Wallis test), but not for Monkey U (p = 0.45, Kruskal-Wallis test). On an individual cell level, a simple linear regression found the same 3/21 cells (Monkey U) and 7/24 cells (Monkey V) were significantly correlated to forthcoming bid as showed significant responses to the presentation of the fractal (see Figure \ref{fig:one_fractal_display}, p < 0.05).}
\label{fig:monkey_bid_fractal_display_quintiles_one}
\end{figure}

When the activity of these populations were binned by the forthcoming bid, only the 7 bid-associated putative dopaminergic cells recorded from Monkey V showed a significant linear relationship between the binned bid and the second component activity to the medium offer fractal (p < 0.0001, adjust $R^2$ = 0.59; Figure \ref{fig:fractal_display_regression_one}A). When normalised to baseline activity (Figure \ref{fig:fractal_display_regression_one}B) this relationship remains significant (p < 0.01, adjusted $R^2$ = 0.35). The 3 putative dopaminergic cells recorded from Monkey U that showed significant bid-related activity do not show a significant linear relationship with the next offer (p = 0.97, adjusted $R^2$ = 0) for the raw counts of cell impulses, but do (p < 0.005, $R^2$ = 0.32) when normalised by baseline activity.

This is a weaker relationship than for when all three fractals are offered (and cells are only grouped by bid, not the magnitude of juice on offer)- see Table \ref{tab:fd_regression_results}. This is likely at least partly due to the reduction in the number of cells and trials recorded, but also potentially the compressed range of bids on these trials compared to the full range of bids made across the three offer magnitudes (for example, compare Figures \ref{fig:monkey_bid_correlation}A) and \ref{fig:single_fractal_bids}A).

\begin{figure}[p] 
\centering    
\includesvg[width=0.945\textwidth,inkscapelatex=false]{Chapter4/Figs/redone_figs/one_fractal_display_regression_fig}
\caption[Putative dopaminergic activity is correlated with forthcoming bid in sessions where only one reward magnitude is offered.]{\textbf{Putative dopaminergic activity is correlated with forthcoming bid in sessions where only one reward magnitude is offered.} Activity for putative dopaminergic cells during single reward level sessions (Monkey U n = 21; Monkey V n = 24) is grouped by the $25^{\text{th}}$ of the budget the forthcoming monkey bid is placed within. Only groupings with more than 10 trials of data are shown. Average firing rate per cell (A) is shown per monkey (Monkey U, blue; Monkey V, orange) with error bars for standard error of mean. Monkey V shows significant correlation between dopaminergic cell impulses and forthcoming bid (orange, p < 0.0001, adjusted $R^2$ = 0.59) though Monkey U does not (blue, p = 0.97, adjusted $R^2$ = 0) using a simple linear regression model. When impulses are normalised to cell's pre-trial baseline activity in B) both monkeys show significant population correlation (Monkey U, blue, p < 0.005, adjusted $R^2$ = 0.32; Monkey V, orange, p < 0.01, adjusted $R^2$ = 0.35) to bid.}
\label{fig:fractal_display_regression_one}
\end{figure}

To directly test whether the dopamine response to the medium fractal was itself modulated by session context, mirroring the behavioural context effect reported in Chapter 3 (Figure \ref{fig:single_fractal_bids}), we compared the normalised dopamine response to the medium fractal between single-fractal and multi-fractal sessions. Single-fractal and three-fractal sessions necessarily involved different, non-overlapping populations of recorded neurons, since cells could not be tracked across sessions; this comparison is therefore inherently between-cells rather than within-cell, unlike the corresponding behavioural comparison, which introduces a source of between-cell variability unrelated to session context itself and reduces the sensitivity of this test independent of sample size.

A linear mixed-effects model\nomenclature{LME}{Linear Mixed-Effects} (the same model structure as Equation \ref{eq:lme_general}, with session type in place of monkey bid as the fixed-effect predictor; trial-level, medium-fractal trials only) showed no significant effect of session context for either monkey (Table \ref{tab:context_comparison}; Monkey U: $\beta = -0.018 \pm 0.140$, $p = 0.898$; Monkey V: $\beta = -0.055 \pm 0.055$, $p = 0.323$). A non-significant result does not in itself quantify evidence for the absence of an effect, therefore we additionally computed Bayes factors ($B_{10}$) using one-sample Bayesian $t$-tests on per-cell mean normalised responses (\cite{rouder2009bayesian}). Both monkeys showed evidence favouring the null of no context effect (Monkey U: $B_{10} = 0.212$, $B_{01} = 4.72$; Monkey V: $B_{10} = 0.308$, $B_{01} = 3.25$), exceeding the conventional threshold of $B_{01} > 3$ for \textit{some evidence} in \cite{jeffreys1961theory}.

\begin{table}[htbp]
\centering
\begin{tabular}{lcccc}
\hline
\textbf{Monkey} & \textbf{LME Estimate} & \textbf{LME $p$-value} & \textbf{$B_{10}$} & \textbf{$B_{01}$} \\
\hline
Monkey U & $-0.018 \pm 0.140$ & $0.898$ & $0.212$ & $4.72$ \\
Monkey V & $-0.055 \pm 0.055$ & $0.323$ & $0.308$ & $3.25$ \\
\hline
\end{tabular}
\caption[No effect of session context on dopamine response to the medium fractal.]{\textbf{No effect of session context on dopamine response to the medium fractal.} Linear mixed-effects model estimates (the same structure as Equation \ref{eq:lme_general}, with session type in place of monkey bid) and Bayes factors ($B_{10}$, $B_{01}$; one-sample Bayesian t-tests on per-cell mean normalised responses, Cauchy prior $r=\sqrt{2}/2$; \cite{rouder2009bayesian}) comparing the dopamine response to the medium fractal between single-fractal and multi-fractal sessions. Evidence labels follow the convention of \cite{jeffreys1961theory}.}
\label{tab:context_comparison}
\end{table}

Finally, to confirm that the dopamine response to the medium fractal tracked offer value at all within single-fractal sessions, we fitted the same mixed-effects model as in Equation \ref{eq:lme_general}, using the same cell population as Figure \ref{fig:fractal_display_regression_one} (Monkey U n = 21 cells; Monkey V n = 24 cells). This showed a significant positive relationship between bid and normalised firing rate for Monkey V ($\beta = 0.584 \pm 0.253$, $p = 0.021$), and a positive but non-significant trend for Monkey U ($\beta = 0.189 \pm 0.109$, $p = 0.084$; Table \ref{tab:medium_fractal_regression}).

\begin{table}[htbp]
\centering
\begin{tabular}{lcccc}
\hline
\textbf{Monkey} & \textbf{$n$ cells} & \textbf{$n$ trials} & \textbf{Slope} & \textbf{$p$-value} \\
\hline
Monkey U & $21$ & $3589$ & $0.189 \pm 0.109$ & $0.084$ \\
Monkey V & $24$ & $841$ & $0.584 \pm 0.253$ & $0.021$ \\
\hline
\end{tabular}
\caption[Trial-level regression of dopamine response on bid, medium fractal, single-fractal sessions.]{\textbf{Trial-level regression of dopamine response on bid, medium fractal, single-fractal sessions.} Linear mixed-effects model estimates as in Equation \ref{eq:lme_general} for putative dopaminergic neurons recorded during single-fractal sessions.}
\label{tab:medium_fractal_regression}
\end{table}

Taken together, these results indicate that the dopamine response to the medium fractal tracked offer value, at least for Monkey V, but did not show the context-dependent modulation seen behaviourally in Monkey V's bids (Chapter 3, Figure \ref{fig:single_fractal_bids}, $d = 0.467$). This suggests the dopamine signal encodes offer value in a manner that is robust to, rather than adapted to, the surrounding session context — an absolute rather than context-relative signal, at least for the medium fractal under the conditions tested here.

\subsubsection{Dopaminergic Responses Split by Offer Value}

To test these hypotheses, we also performed the same linear modelling as in Figures \ref{fig:fractal_display_regression} and \ref{fig:fractal_display_regression_one} for sessions when all three offer magnitudes are available to the monkey, but split the regression by the fractal shown (Figure \ref{fig:fractal_display_regression_split}). 

\begin{figure}[p] 
\centering    
\includesvg[width=0.93\textwidth,inkscapelatex=false]{Chapter4/Figs/redone_figs/fd_regression_by_fractal_fig}
\caption[Putative dopaminergic activity at offer epoch is correlated with monkey's bid for each offered reward magnitude.]{\textbf{Putative dopaminergic activity at offer epoch is correlated with monkey's bid for each offered reward magnitude.} When split by reward (highest top, lowest bottom) offered, monkey's bids remain correlated with normalised activity of bid-responsive putative dopaminergic neurons. For Monkey U (left, blue; n = 32 neurons) responses to the display of the high (p < 0.01, adjusted $R^2$ = 0.52) and medium (p < 0.0005, adjusted $R^2$ = 0.80) offer fractals were significant. Responses to the low fractal offer were not significantly correlated to forthcoming bid (p = 0.15, adjusted $R^2$ = 0.15). For Monkey V (n = 41 neurons), dopaminergic responses to all three offers were significantly correlated to with the forthcoming bid (low, p < 0.005, adjusted $R^2$ = 0.62; medium, p < 0.005, adjusted $R^2$ = 0.69; high, p < 0.005, adjusted $R^2$ = 0.71). All plots show mean normalised firing rate +/- standard error of the mean.}
\label{fig:fractal_display_regression_split}
\end{figure}

Due to a reduction in the number of bids made at various percentiles of the total bidding space per fractal, we reduced the number of bins from 25 in Figure \ref{fig:fractal_display_regression} to 20 per fractal (each representing a range of 0.06ml of water. The same 32 (Monkey U) and 41 (Monkey V) cells are used whose activity was significantly correlated with monkey bids (regardless of fractal value). We only take into account normalized second-component activity per cell in a simple linear regression of the form shown in Equation \ref{eq:normalized_firing_rate_regression_fractal}:

\begin{align}
\frac{\text{Firing Rate}}{\text{Baseline Activity}} | \text{Fractal Offer} &\sim \beta_0 + \beta_1(\text{Monkey Bid Bin}) + \varepsilon \label{eq:normalized_firing_rate_regression_fractal}
\end{align}

For Monkey V, all three linear regressions show a significant correlation between the forthcoming bid of the monkey for a trial and the impulses of the second component of putative dopaminergic activity (low: p < 0.005, adjusted $R^2$ = 0.36; medium: p < 0.05, adjusted $R^2$ = 0.18; high: p < 0.001, adjusted $R^2$ = 0.62). For Monkey U, the relationship for trials when the medium (p < 0.0001, adjusted $R^2$ = 0.67) and highest (p < 0.01, adjusted $R^2$ = 0.31) was significant, but there was not a significant relationship between the lowest juice offer and the subsequent bid (p = 0.14, adjusted $R^2$ = 0.07). The full results of these regressions are shown in Table \ref{tab:fd_regression_split_fractal}.

\begin{table}[htbp]
\centering
\begin{tabular}{llcccc}
\hline
\textbf{Monkey }& \textbf{Offer Value} & \textbf{Intercept} & \textbf{Slope} & \textbf{Adj} $R^2$ & $p$\textbf{-value} \\
\hline
Monkey U & High & $-0.021 \pm 0.306$ & $1.488 \pm 0.498$ & 0.306 & $8.26 \times 10^{-3}$ \\
Monkey U & Medium & $-0.707 \pm 0.109$ & $1.147 \pm 0.182$ & 0.671 & $6.14 \times 10^{-6}$ \\
Monkey U & Low & $-0.594 \pm 0.249$ & $0.639 \pm 0.416$ & 0.067 & $1.42 \times 10^{-1}$ \\
\hline
Monkey V & High & $-1.831 \pm 0.534$ & $3.593 \pm 0.764$ & 0.619 & $5.13 \times 10^{-4}$ \\
Monkey V & Medium & $-0.522 \pm 0.144$ & $0.552 \pm 0.240$ & 0.183 & $3.40 \times 10^{-2}$ \\
Monkey V & Low & $-0.628 \pm 0.071$ & $0.407 \pm 0.118$ & 0.365 & $2.85 \times 10^{-3}$ \\
\hline
\end{tabular}
\caption[Significant relationship between dopaminergic cell activity and forthcoming monkey bid at each offer value level.]{\textbf{Significant relationship between dopaminergic cell activity and forthcoming monkey bid at each offer value level.} Forthcoming bid for juice tends to be correlated with the normalised second-component activity per monkey and juice magnitude offer. Simple linear regression grouped by monkey and trial good offered. Only the relationship between the monkey's subsequent bid and the dopaminergic cell activity on trials for the lowest fractal offer for Monkey U was not significant (p = 0.14).}
\label{tab:fd_regression_split_fractal}
\end{table}

As with the central claim above, these regressions pool trials within each fractal-value subset as though they were independent observations. We therefore also fitted the same mixed-effects model as in Equation \ref{eq:lme_general}, separately within each fractal-value subset:

\begin{align}
\frac{\text{Firing Rate}}{\text{Baseline Activity}} \Big| \text{Fractal Offer} = \beta_0 + \beta_1(\text{Monkey Bid}) + u_{\text{cell}} + \varepsilon, \qquad u_{\text{cell}} \sim \mathcal{N}(0, \sigma^2_{\text{cell}}) \label{eq:lme_fractal}
\end{align}

This confirmed significant positive relationships for five of the six monkey-by-fractal combinations (Monkey U: low $\beta = 0.560 \pm 0.277$, $p = 0.044$; medium $\beta = 1.06 \pm 0.235$, $p = 8.09 \times 10^{-6}$; high $\beta = 1.79 \pm 0.487$, $p = 2.53 \times 10^{-4}$; Monkey V: low $\beta = 0.392 \pm 0.145$, $p = 0.007$; medium $\beta = 0.470 \pm 0.191$, $p = 0.014$). The exception was Monkey V's response to the high-value fractal, which showed the strongest relationship of the six under the naive regression above (Table \ref{tab:fd_regression_split_fractal}) but is no longer significant once cell-level clustering is accounted for ($\beta = 0.160 \pm 0.578$, $p = 0.782$).

\begin{table}[htbp]
\centering
\begin{tabular}{llccc}
\hline
\textbf{Monkey }& \textbf{Offer Value} & \textbf{Intercept} & \textbf{Slope} & $p$\textbf{-value} \\
\hline
Monkey U & High & $-0.0836 \pm 0.416$ & $1.79 \pm 0.487$ & $2.53 \times 10^{-4}$ \\
Monkey U & Medium & $-0.631 \pm 0.149$ & $1.06 \pm 0.235$ & $8.09 \times 10^{-6}$ \\
Monkey U & Low & $-0.606 \pm 0.118$ & $0.560 \pm 0.277$ & $4.41 \times 10^{-2}$ \\
\hline
Monkey V & High & $0.875 \pm 0.527$ & $0.160 \pm 0.578$ & $7.82 \times 10^{-1}$ \\
Monkey V & Medium & $-0.455 \pm 0.146$ & $0.470 \pm 0.191$ & $1.40 \times 10^{-2}$ \\
Monkey V & Low & $-0.619 \pm 0.0774$ & $0.392 \pm 0.145$ & $6.70 \times 10^{-3}$ \\
\hline
\end{tabular}
\caption[Mixed-effects model estimates for dopaminergic cell activity against forthcoming monkey bid at each offer value level.]{\textbf{Mixed-effects model estimates for dopaminergic cell activity against forthcoming monkey bid at each offer value level.} Results from the linear mixed-effects model described in Equation \ref{eq:lme_fractal}, fitted separately per monkey and fractal offer value, with a random intercept per cell. Adj $R^2$ is not directly applicable to this model type and is not included. All conditions show a significant positive relationship except Monkey V's response to the high-value fractal ($p = 0.782$), which was significant under the naive regression in Table \ref{tab:fd_regression_split_fractal} but is not once cell-level clustering is accounted for.}
\label{tab:lme_regression_split_fractal}
\end{table}

These data show that the activity of a subset of dopaminergic cells is not only encoding just the physical properties of the reward (i.e., the volume of juice offered; Figure \ref{fig:fractal_display}), but also the \textit{subjective value} that the monkey places upon it. The number of impulses recorded from these cells is correlated to the monkey's indication of their value of the offer across all trials, \textit{but also} when grouped by the magnitude of the reward on offer.

To verify whether dopaminergic activity was independent of reward magnitude, we also wanted to study whether the inverse is true; when monkeys make similar offers for different goods, do these cells show similar responses?

As the resolution of the bids made was sufficiently detailed (to the thousandth of the bidding space) for them to be treated as continuous, different fractal offers were grouped if they fell within 2.5\% of the total bidding space (a range of 0.03ml of water) as judged to be 'equivalent.' The three combinations of grouping are therefore between trials where the high reward magnitude is shown and the low reward magnitude is shown ('H-L', red, left-most), where the high and medium reward magnitudes are shown ('H-M', orange, middle), and where the medium and low reward magnitudes are shown ('M-L', yellow, right-most). The second component activity of bid-associated (see Figure \ref{fig:fractal_display_regression_split}) dopaminergic cells per trial is normalised and the difference between activity is calculated. We subtract the activity of the trial with the lower reward offer from the activity on the trial with the greater reward offer, i.e., in 'H-L', the normalised dopaminergic activity on the trial where the lowest magnitude reward fractal is shown is subtract from the one where the highest magnitude reward fractal is offered.

For both monkeys, a Wilcoxon signed-rank test did not show significant differences for any of the pairings (p > 0.05; Figure \ref{fig:fd_similar_bids_activity}). A non-significant result establishes that firing rate differences between matched-bid trial pairs were not statistically reliable, but does not in itself quantify positive evidence for the absence of a difference. To address this, we computed Bayes factors (BFs\nomenclature{BF}{Bayes Factor}) using one-sample Bayesian $t$-tests on per-cell mean firing rate differences, with a default Cauchy prior ($r = \sqrt{2}/2$; \citealt{rouder2009bayesian}). This approach allowed us to ask not merely whether a difference could be ruled out at a significance threshold, but how strongly the data favoured the null hypothesis of equivalent firing relative to the alternative.

BFs ranged from 0.785 to 3.280 across the six monkey-comparison combinations (Table \ref{tab:bayes_factor_matched_bids}), with five of six below, and one at the boundary of, the conventional threshold of 3 for \textit{some evidence} of a difference\footnote{Following \citet{jeffreys1961theory}: BF $>$ 3 = some evidence, BF $>$ 10 = strong evidence, BF $>$ 30 = very strong evidence for one hypothesis over another.}. The results are therefore genuinely equivocal: the data neither strongly support nor decisively rule out a residual magnitude signal. Importantly, the fractals used here differ substantially in objective reward magnitude, and the monkey's mean bid for each fractal differs accordingly (Figure \ref{fig:monkey_bid_correlation}). Under a pure magnitude-encoding hypothesis, one would expect firing rate differences to persist even when individual bids are matched, and the resulting BFs to be large. The fact that no BF substantially exceeded 3 across all comparisons, despite these large objective differences between conditions, is consistent with neurons tracking subjective value as expressed through the bid rather than reward magnitude per se. On the display of the good offered, dopaminergic cells encode the bid to be made which, given the incentive-compatible nature of the BDM is equivalent to the subjective value of the offer.

\begin{figure}[htbp!] 
\centering    
\includegraphics[width=1.0\textwidth]{Chapter4/Figs/redone_figs/fd_similar_bids_fig.png}
\caption[Putative dopaminergic neurons exhibit similar activity for similar bids across offered reward magnitudes.]{\textbf{Putative dopaminergic neurons exhibit similar activity for similar bids across offered reward magnitudes.} When trials in which similar bids are made for different fractal offers are compared, there is no significant difference in responses of putative dopaminergic neurons. Differences in impulses in the second component of the dopaminergic neuron response are plotted per similar trial-pair for Monkey U (n = 32 neurons, left) and Monkey V (n = 41 neurons, right) across all 3 difference offer comparisons of high (H), medium (M) and low (L) fractal offer. For each comparison, the responses to the higher reward magnitude offer is taken as the reference. No significant differences (p > 0.05, Wilcoxon signed-rank test) were found over the population of the trials for each fractal-fractal comparison.}
\label{fig:fd_similar_bids_activity}
\end{figure}

\begin{table}[h]
\centering
\begin{tabular}{llcc}
\hline
\textbf{Monkey} & \textbf{Comparison} & \textbf{Bayes Factor} & \textbf{Evidence} \\
\hline
Monkey U & H-L & 3.280 & some evidence ($H_1$) \\
Monkey U & H-M & 0.785 & --- \\
Monkey U & M-L & 1.610 & --- \\
\hline
Monkey V & H-L & 1.440 & --- \\
Monkey V & H-M & 0.961 & --- \\
Monkey V & M-L & 0.793 & --- \\
\hline
\end{tabular}
\caption[Bayes factors for matched-bid firing rate comparisons.]{Bayes factors ($B_{10}$) from one-sample Bayesian $t$-tests on per-cell mean firing rate differences for matched-bid trial pairs, computed using a Cauchy prior on effect size ($r = \sqrt{2}/2$; \citealt{rouder2009bayesian}). Evidence labels follow \citet{jeffreys1961theory}. H = high, M = medium, L = low reward magnitude fractals.}
\label{tab:bayes_factor_matched_bids}
\end{table}

\clearpage
\section{Dopaminergic Responses to Later Tasks Epochs}

Although our goal was to study the relationship between dopaminergic neuron activity and the monkey's subjective valuation (bid) of an offered good, we also present analyses of later task epochs. These analyses demonstrate both canonical dopaminergic responses throughout the task and verify the monkey's understanding of the information presented.

\subsection{Dopaminergic neurons encode profit at the trial outcome}

At the reveal of the trial outcome, the monkey already knows the offered reward magnitude and their own bid. They then learn the final piece of information, the computer's random bid, which determines whether they receive juice and how much water they forgo. This is shown to the monkey as a green bar that crosses the hashed bidding space above or below their own purple bidding indicator bar (see Figure \ref{fig:task_epochs}). If the monkey's purple bar is above (i.e., the monkey made a higher bid than the computer), the pattern on the bidding space underneath the second-price computer bid is reversed, indicating the amount of water the monkey 'pays' in order to receive their juice reward. As the computer bid is drawn randomly, this represents the last piece of unpredicted information in the task.

In Figure \ref{fig:winlose_reveal}A), responses to this outcome reveal are grouped by whether or not the monkey won the auction. In any willingness-to-pay auction, a participant can only stand to make a profit when they win the offered good\footnote{After a losing bid, the participant gets to keep their endowment as was. This slightly ignores the potential time and effort discounting of taking part in an auction trial, and also that monkeys only receive their 'budget' at the start of the trial; see Figure \ref{fig:winlose_reveal_by_fractal}B).}. Both monkeys show a clear increase in normalised population response to the reveal of the trial outcome on trials where they made the higher bid and won the offered juice reward (Wilcoxon signed-rank test, p < 0.0001 for both monkeys). Interestingly, the population of cells recorded from Monkey U showed a decrease in activity when losing a trial (Wilcoxon signed-rank test, p < 0.005), whereas the population of recorded Monkey V neurons showed the opposite- a clear, modest increase in the number of impulses recorded (p < 0.0001). 

Both monkeys showed highly significant changes in putative dopaminergic cell activity in trials when they won the auction compared to when they lost (Wilcoxon signed-rank test, p < 0.0001 for both). 72/123 putative dopaminergic neurons recorded from Monkey U and 46/145 from Monkey V reached individual significance for a change in activity under the same test.

\begin{figure}[p] 
\centering    
\includesvg[width=0.92\textwidth,inkscapelatex=false]{Chapter4/Figs/redone_figs/winlose_fig}
\caption[Putative dopaminergic neurons show greater activation when monkeys win their auction trial.]{\textbf{Putative dopaminergic neurons show greater activation when monkeys win their auction trial.} A) Normalised activity of dopaminergic neurons (Monkey U n = 123; Monkey V n = 145) shows an increase when the monkey is revealed to have won a trial (purple, p < 0.0001 for both monkeys; Wilcoxon signed-rank test). This increase is especially marked compared to the activity of the same neurons when a loss is revealed. In losing trials, the neurons recorded from Monkey U show a significant decrease in activity (black; p < 0.005 Wilcoxon signed-rank test), whereas the same test for Monkey V shows a significant (p < 0.0001) increase in activity. The same data for losing trials is shown in black in B) while the winning trials are split by the amount of 'profit' a monkey makes on the trial. A significant difference between the large and small profit outcomes was revealed on a post hoc Dunn test for both monkeys (p < 0.0001).}
\label{fig:winlose_reveal}
\end{figure}

The averaged trial responses of those 72 and 46 neurons are shown for both trials in which the monkey wins the auction (top) and loses the auction (bottom) in Figure \ref{fig:winlose_reveal_raster}. The same averaged neuronal activity data is shown in Figure \ref{fig:winlose_profit_raster}, split by the profit the monkey made on those trials, as described above.

\begin{figure}[p] 
\centering    
\includegraphics[width=1.0\textwidth]{Chapter4/Figs/redone_figs/win_raster_fig}
\caption[Averaged responses of putative dopaminergic neurons with significant responses to the reveal of trial outcome.]{\textbf{Averaged responses of putative dopaminergic neurons with significant responses to the reveal of trial outcome.} A Wilcoxon signed-rank test revealed that the activity of n = 72/123 neurons (Monkey U) and n = 46/145 neurons (Monkey V) was significantly affected by the reveal of the trial outcome. The responses for each of these individual neurons averaged over all trials is shown per row, split by whether the monkey won the trial (top) or lost (bottom).}
\label{fig:winlose_reveal_raster}
\end{figure}

\begin{figure}[p] 
\centering    
\includegraphics[width=1.0\textwidth]{Chapter4/Figs/redone_figs/profit_raster_fig}
\caption[Averaged responses of putative dopaminergic neurons with significant responses to the reveal of trial outcome split by profit.]{\textbf{Averaged responses of putative dopaminergic neurons with significant responses to the reveal of trial outcome split by profit.} As in Figure \ref{fig:winlose_reveal_raster} neurons were selected by significant changes in activity under a Wilcoxon signed-rank test when the trial outcome was revealed. Here split by the monkey's profit on each trial as displayed in Figure \ref{fig:winlose_reveal}B). Responses to trials where the monkey made a profit of >33\% of the bidding space is shown at the top, a small profit of <= 33\% of the bidding space is shown in the middle, and the same loss data as in Figure \ref{fig:winlose_reveal_raster} bottom is also shown.}
\label{fig:winlose_profit_raster}
\end{figure}

For two fully random bidders, the expected difference between bids given that a defined bid is higher than the other is 33\% of the budget. Therefore, we split the population activity into 'lose', 'small profit' (a profit of less than 33\% of the total bidding space), and 'large profit' (a profit of more than 33\% of the total bidding space; Figure \ref{fig:winlose_reveal}B). A Kruskal-Wallis test assuming that dopaminergic cells would show greater activity to a large profit over a small profit and both over losing a trial was significant for both monkeys (p < 0.0001). Visual inspection of \ref{fig:winlose_reveal}B) showed the pattern we expected above for Monkey U after the reveal of the outcome of the trial. Monkey V however had similar response patterns for the 'small profit' and 'loss' conditions. A post hoc Dunn test confirmed that for Monkey U, all pairwise comparisons were significant (p < 0.0001). For Monkey V, while the large profit condition produced significantly greater activity than both the small profit and the loss conditions (p < 0.0001), there was no significant difference between small profit and the loss conditions (p = 0.1037).

To investigate this, we analysed the behaviour of recorded putative dopaminergic cells from both monkeys broken down by the reward magnitude on offer each trial. Figure \ref{fig:winlose_reveal_by_fractal}A) shows the normalised impulses of these cells after the monkey is revealed to have won the auction in a BDM trial. Unsurprisingly, the population activity for each offer level shows an increase in firing rate (Wilcoxon signed-rank test, Monkey V, low reward offer p < 0.0005, all others p < 0.0001); for any fractal, if the monkey has bid their subjective value and won the trial, they have made a profit. A Kruskal-Wallis test on the rank order of dopaminergic cell activity was significant for both monkeys (p < 0.0001) with a post hoc Dunn test showing significant differences between the highest and medium reward magnitude offer and the lowest magnitude offer (p < 0.0001), though this test was not significant for differences between the highest and medium reward fractals for either monkey (Monkey U p = 0.183; Monkey V p = 0.0672).

For trials in which the monkey bids lower than the computer and loses the auction, the reverse is true (Figure \ref{fig:winlose_reveal_by_fractal}B). For Monkey U all three offer values provoked significant changes in firing rate from the pre-reveal baseline (low value: increase, p < 0.0005, else: decrease, p < 0.0001). In contrast, for Monkey V only the lowest and medium offer fractals led to an increase (p < 0.0001) in firing rate, whereas these was no significant effect when the monkey was revealed to have lost a trial for the highest reward offer (p = 0.830; all Wilcoxon signed-rank test).

\begin{figure}[p] 
\centering    
\includesvg[width=0.86\textwidth,inkscapelatex=false]{Chapter4/Figs/redone_figs/winlose_fractal_fig}
\caption[Putative dopaminergic neurons activity at task outcome reveal shows ordering by reward offered.]{\textbf{Putative dopaminergic neurons activity at task outcome reveal shows ordering by reward offered.} A) When monkeys (Monkey U, right; Monkey V, left) win the auction a clear increase in putative dopaminergic neuron activity is seen (all conditions p < 0.0005). For both monkeys the ordering of normalised dopaminergic activity was significant (Kruskal-Wallis test; p < 0.0001) and a post-hoc Dunn test showed significant differences in activity between the highest-lowest and medium-lowest reward levels (p < 0.0001) though not between the highest and medium reward levels (Monkey U p = 0.183, Monkey V p = 0.0672). B) When losing a trial and forgoing the highest and medium reward levels, the normalised activity recorded from Monkey U showed significant decreases (Wilcoxon signed-rank test p < 0.0001). When Monkey U lost vs. the lowest reward level, these cells showed a significant \textit{increase} in activity (p < 0.0005). The same tests for recordings from Monkey V showed significant increases in firing rate when losing to the computer's bid on trials where the lowest two reward levels were shown (p < 0.0001) but no significant change in activity (p = 0.830) when the monkey did not win when offered the highest potential juice reward.}
\label{fig:winlose_reveal_by_fractal}
\end{figure}

A Kruskal-Wallis test between groups was significant for both monkeys (Monkey U p < 0.0001, Monkey V p < 0.0005) though a post hoc Dunn test showed that for Monkey U there was no significant difference between the responses to trials where the monkey lost an auction for the highest and medium reward magnitude fractals (p = 0.607). All other comparisons were significant (Monkey V high - medium comparison p < 0.01, else p < 0.0001).

A possible explanation for the behaviour of neurons observed in loss trials for Monkey V (Figure \ref{fig:winlose_reveal}, to a lesser extent for Monkey U Figure \ref{fig:winlose_reveal_by_fractal}B) could be that for the lowest offer trials, the monkey is less interested in the (possible profit from the) reward of winning a trial, and more focused on the potential greater payoff on later trials. The win lose reveal comes after a substantial 4.5s (Figure \ref{fig:task_epochs}) bidding phase (during a long 12s total length trial) during which they are unable to win any extra liquid to sate their thirst. The positive 'reward' prediction error may be at least partly explained by the monkey unable to accurately predict the end of this epoch and encouraged that they can now receive (when losing) the value of their total budget. In addition, they are closer to the onset of a new trial, which may yield the much more rewarding highest magnitude juice offer for them to bid for (see Figure \ref{fig:fixation_cross} and the increase in dopaminergic firing rate at the beginning of a new trial).

This may be supported by Figure \ref{fig:winlose_loss_regression}A) in which the responses of putative dopaminergic neurons are plotted given the tercile of the monkey's bid (in thirds of the total bidding space) on trials where the monkey loses. For both monkeys, an increase in activity is only seen at the outcome reveal when the monkey makes a low bid for the offered good (Monkey U p < 0.0005, Monkey V p < 0.0001, Wilcoxon signed-rank test) with the tercile of the monkey's bid significant on a Kruskal-Wallis test (Monkey U p < 0.0001, Monkey V p < 0.0005). The implied low subjective value of these trials (due to whichever cause) is almost complete, and the monkey can hope to receive a higher-valued trial next.

A simple linear regression of the form in Equation \ref{eq:normalized_firing_rate_bid} of the binned monkey bid was significant for both monkeys (Monkey U p < 0.0001, adjusted $R^2$ = 0.783; Monkey V p < 0.0005, adjusted $R^2$ = 0.418).

\begin{align}
\frac{\text{Firing Rate}}{\text{Baseline Activity}} | \text{Lose} &\sim \beta_0 + \beta_1(\text{Monkey Bid Bin}) + \varepsilon \label{eq:normalized_firing_rate_bid}
\end{align}

\begin{figure}[p] 
\centering    
\includesvg[width=0.9\textwidth,inkscapelatex=false]{Chapter4/Figs/redone_figs/loss_regression}
\caption[Putative dopaminergic neurons show greater decreases in firing rate when monkeys fail to win a highly valued good.]{\textbf{Putative dopaminergic neurons show greater decreases in firing rate when monkeys fail to win a highly valued good.} A) Normalised population activity of cells which showed significant changes in activity upon the reveal of the win/lose stimulus are plotted, coloured by the tercile of the bidding space in which monkeys placed their bid. For both monkeys, a decrease in activity is shown when the trial is lost but the monkey placed a bid in the upper $2/3^{\text{rds}}$ of the bidding space (teal, yellow). For both monkeys an increase in activity was shown when losing having made a low bid (dark blue; Monkey U p < 0.0005, Monkey V p < 0.0001; Wilcoxon signed-rank test). A Kruskal-Wallis test for ordering of activity was significant for both (Monkey U p < 0.0001, Monkey V p < 0.0005). B) The total bidding space is split into $25^{\text{ths}}$ and the monkeys bid on a trial binned (Monkey U blue, left; Monkey V orange, right). For both monkeys the regression in Equation \ref{eq:normalized_firing_rate_bid} is highly significant (p < 0.0001) with an adjusted $R^2$ of 0.783 for Monkey U and 0.418 for Monkey V.}
\label{fig:winlose_loss_regression}
\end{figure}

Of course, it is also possible that monkeys may not be bidding in a fully incentive-compatible manner and be relieved at the loss of a trial on which they may have overbid for the offered good.

To further study how putative dopaminergic neurons react to the outcome of a BDM trial, we performed two regressions of the normalised firing rate: against the profit the monkey made, and the computer's random bid, on any trial (Equations \ref{eq:normalized_firing_rate_profit} and \ref{eq:normalized_firing_rate_cb}):

\begin{align}
\frac{\text{Firing Rate}}{\text{Baseline Activity}} | \text{Win} &\sim \beta_0 + \beta_1(\text{Monkey Profit Bin}) + \varepsilon \label{eq:normalized_firing_rate_profit}\\
\frac{\text{Firing Rate}}{\text{Baseline Activity}} | \text{Lose} &\sim \beta_0 + \beta_1(\text{Computer Bid Bin}) + \varepsilon \label{eq:normalized_firing_rate_cb}
\end{align}

The variance in neuronal activity at this epoch when the monkey makes a winning bid is very well explained by the profit the monkey makes on the trial\footnote{Shown as a fraction of the maximum possible profit, i.e., when the computer bid is zero} (Monkey U p < 0.0001, adjusted $R^2$ = 0.923; Monkey V p < 0.0001, adjusted $R^2$ = 0.913). In contrast, when the monkey loses the auction the offer made by the computer in a trial does not have a significant effect on the responses from either monkey (Monkey U p = 0.717, adjusted $R^2$ = -0.043; Monkey V p = 0.125, adjusted $R^2$ = 0.082). The neurons are acting as if the monkey understands the mechanics of the task and how much profit they are about to receive on a trial, or that it doesn't matter by how much they lose an auction.

\begin{figure}[htbp!] 
\centering    
\includesvg[width=1.0\textwidth,inkscapelatex=false]{Chapter4/Figs/redone_figs/winlose_regression_fig}
\caption[Putative dopaminergic activity is strongly correlated to the monkey's profit on a BDM trial.]{\textbf{Putative dopaminergic activity is strongly correlated to the monkey's profit on a BDM trial.} In A) the binned profit (monkey's bid minus computer bid) is split into $25^{\text{ths}}$ of the total bidding space. This binned profit that the monkey makes on a BDM trial shows a strong, significant correlated with the normalised firing rate of putative dopaminergic neurons (both monkeys p < 0.0001). In contrast, in trials where the monkey loses, the value of the higher computer bid has no significant relationship with the activity of these cells (Monkey U p = 0.717, Monkey V p = 0.125).}
\label{fig:winlose_profit_regression}
\end{figure}

The full results of the three linear regressions for Equations \ref{eq:normalized_firing_rate_bid}, \ref{eq:normalized_firing_rate_profit}, \ref{eq:normalized_firing_rate_cb} (Figures \ref{fig:winlose_loss_regression} and \ref{fig:winlose_profit_regression}) are shown in Table \ref{tab:wl_regression_results}.

\begin{table}[htbp]
\centering
\begin{tabular}{lllccc}
\hline
\textbf{Monkey} & \textbf{Ind. Variable} & \textbf{Intercept} & \textbf{Slope} & \textbf{Adj} $R^2$ & $p$\textbf{-value} \\
\hline
Monkey U & Monkey Bid & $0.27 \pm 0.04$ & $-0.71 \pm 0.08$ & 0.783 & $1.25 \times 10^{-8}$ \\
Monkey U & Monkey Profit & $0 \pm 0.05$ & $1.53 \pm 0.11$ & 0.923 & $1.52 \times 10^{-10}$ \\
Monkey U & Computer Bid & $-0.05 \pm 0.13$ & $0.08 \pm 0.21$ & -0.043 & $7.17 \times 10^{-1}$ \\
\hline
Monkey V & Monkey Bid & $0.39 \pm 0.07$ & $-0.52 \pm 0.12$ & 0.418 & $3.81 \times 10^{-4}$ \\
Monkey V & Monkey Profit & $0.02 \pm 0.04$ & $1.04 \pm 0.07$ & 0.913 & $1.01 \times 10^{-11}$ \\
Monkey V & Computer Bid & $0.48 \pm 0.13$ & $-0.31 \pm 0.19$ & 0.082 & $1.25 \times 10^{-1}$ \\
\hline
\end{tabular}
\caption[Neuronal response at the win/lose reveal is only dependent on monkey's subjective value.]{\textbf{Neuronal response at the win/lose reveal is only dependent on monkey's subjective value.} Full statistics for the regressions performed against the neuronal response of each monkey to the reveal of the outcome of the trial. Only variables related to the monkey's subjective valuation of the offered good (their bid, or the profit made) are significantly correlated with changes in neuronal firing when the monkey either wins ('Monkey Profit') or loses (else) the trial. See Equations \ref{eq:normalized_firing_rate_bid}, \ref{eq:normalized_firing_rate_profit}, \ref{eq:normalized_firing_rate_cb}.}
\label{tab:wl_regression_results}
\end{table}

We also fitted mixed-effects equivalents of Equations \ref{eq:normalized_firing_rate_bid}, \ref{eq:normalized_firing_rate_profit}, and \ref{eq:normalized_firing_rate_cb}, following the same general structure as Equation \ref{eq:lme_general}:

\begin{align}
\frac{\text{Firing Rate}}{\text{Baseline Activity}} \Big| \text{Outcome} = \beta_0 + \beta_1(\text{Predictor}) + u_{\text{cell}} + \varepsilon, \qquad u_{\text{cell}} \sim \mathcal{N}(0, \sigma^2_{\text{cell}}) \label{eq:lme_winlose}
\end{align}

Outcome is either Win or Lose, and Predictor is the relevant subjective-value variable for that trial type (monkey bid, monkey profit, or computer bid).

For trials where the monkey lost, this confirmed a significant negative relationship between normalised firing rate and monkey bid for both monkeys ($\beta = -0.658 \pm 0.118$, $p = 2.45 \times 10^{-8}$ for Monkey U; $\beta = -0.433 \pm 0.103$, $p = 2.60 \times 10^{-5}$ for Monkey V), consistent with the naive regression above. For the computer's bid on loss trials, however, the mixed-effects model diverges from the naive result: Monkey U remained non-significant ($\beta = -0.161 \pm 0.102$, $p = 0.116$), but Monkey V showed a significant negative relationship ($\beta = -0.590 \pm 0.119$, $p = 7.60 \times 10^{-7}$) that the naive population-level regression above (Table \ref{tab:wl_regression_results}) did not detect.

This is not what a purely payout-driven account would predict as on loss trials the monkey receives the same full budget regardless of the computer's bid magnitude. There is no objective reward-based reason for putative dopaminergic neurons to track it at all. One candidate explanation is a missed-opportunity signal: a low computer bid on a lost trial means the auction could have been won cheaply, whereas a high computer bid means the loss was harder to avoid regardless of strategy. Humans show a comparable pattern of anticipated loser regret in auction settings, from losing at an affordable price (\cite{filiz-ozbayAuctionsAnticipatedRegret2007}), and engage measurably in strategic reasoning about a competitor's bid in second-price auctions specifically (\cite{gretherMentalProcessesStrategic2007}). While bidding one's true utility should preclude formal regret, random utility theory holds that the value placed on a choice is sampled from a distribution at each evaluation and is inherently noisy (Appendix A, Section \ref{sec:rum}; \cite{mcfaddenEconometricModelsProbabilistic1982}). If a monkey underbids their true value, a low computer bid is precisely the affordable missed profit scenario that the regret framework introduced in Chapter 2 (Figure \ref{fig:bdm_outcomes}) describes. This also connects to that chapter's broader point that second-price auctions of this kind are notoriously difficult even for humans to grasp strategically (\cite{casonMisconceptionsGameForm2014}; \cite{hamukwalaDesignFactorsInfluencing2019}).

For trials where the monkey won, this confirmed a significant positive relationship between normalised firing rate and profit for both monkeys ($\beta = 1.71 \pm 0.134$, $p = 2.18 \times 10^{-36}$ for Monkey U; $\beta = 1.02 \pm 0.0817$, $p = 8.17 \times 10^{-35}$ for Monkey V), consistent with the naive regression above and remaining the strongest relationship of the three win/lose regressions under both approaches.

\begin{table}[htbp]
\centering
\begin{tabular}{lllcc}
\hline
\textbf{Monkey} & \textbf{Ind. Variable} & \textbf{Intercept} & \textbf{Slope} & $p$\textbf{-value} \\
\hline
Monkey U & Monkey Bid & $0.261 \pm 0.0691$ & $-0.658 \pm 0.118$ & $2.45 \times 10^{-8}$ \\
Monkey U & Monkey Profit & $0.00347 \pm 0.0830$ & $1.71 \pm 0.134$ & $2.18 \times 10^{-36}$ \\
Monkey U & Computer Bid & $0.109 \pm 0.0869$ & $-0.161 \pm 0.102$ & $1.16 \times 10^{-1}$ \\
\hline
Monkey V & Monkey Bid & $0.293 \pm 0.0847$ & $-0.433 \pm 0.103$ & $2.60 \times 10^{-5}$ \\
Monkey V & Monkey Profit & $0.0142 \pm 0.0732$ & $1.02 \pm 0.0817$ & $8.17 \times 10^{-35}$ \\
Monkey V & Computer Bid & $0.506 \pm 0.111$ & $-0.590 \pm 0.119$ & $7.60 \times 10^{-7}$ \\
\hline
\end{tabular}
\caption[Mixed-effects model estimates for neuronal response at the win/lose reveal.]{\textbf{Mixed-effects model estimates for neuronal response at the win/lose reveal.} Results from the linear mixed-effects model described in Equation \ref{eq:lme_winlose}, fitted separately per monkey and predictor, with a random intercept per cell. c.f. Table \ref{tab:wl_regression_results} for the naive regression equivalents.}
\label{tab:lme_wl_regression_results}
\end{table}

\subsection{Dopaminergic cells show significant responses to reward receipt}

The final two epochs of the task (Figure \ref{fig:task_epochs}) are the payouts, first of the juice (if the monkey won the trial), then of any remaining water budget. Both the timing (1.5s delay before the receipt of either) and the liquid reward (given the reveal of the computer bid) to the monkey can be completely predicted. Nevertheless, we might expect a positive reward prediction error as the monkey is unable to keep perfect time and there is always some risk that perhaps the trial will not be paid out\footnote{E.g., maybe the task has changed on that trial, or even that the sugary juice has stuck the liquid taps shut.}. For both of these epochs, there is a visual stimulus (the fractal disappears from the task screen as the juice is delivered; the bidding space and bidding bars disappear as the budget water is delivered) that occurs simultaneously as the relevant liquid tap opens and the monkey receives either juice or water.

In Figure \ref{fig:juice_delivery} we see a clear and significant population response (Wilcoxon signed-rank test; Monkey U p < 0.005, Monkey V p < 0.0001) to removal of the fractal from the task screen and payment of the juice reward (purple) or not (black) depending on whether the monkey won the auction trial. Individual Wilcoxon signed-rank tests yielded 28/123 (Monkey U) and 47/145 (Monkey V) cells with significant increases in activity upon payment of the juice reward\footnote{In this and further figures in this chapter, responses are aligned to the time relative to the visual indication of liquid delivery rather than the liquid delivery \textit{per se}. The time between these two stimuli is minor and of a constant duration across trials.}. Averaged responses over all trials for these cells are shown in Figure \ref{fig:juice_delivery_raster}.

\begin{figure}[p] 
\centering    
\includesvg[width=1.0\textwidth,inkscapelatex=false]{Chapter4/Figs/redone_figs/juice_delivery_fig}
\caption[Putative dopaminergic population activity increases upon payment of the auctioned trial reward juice.]{\textbf{Putative dopaminergic population activity increases upon payment of the auctioned trial reward juice.} For each monkey, the population activity of putative dopaminergic cells shows a clear increase shortly after the payment of the reward juice in trials where the monkey won the auction (purple) compared to trials where the monkey's bid was lower than that of the computer (black). A Wilcoxon signed-rank test between these two conditions showed significance for both monkeys (Monkey U p < 0.005, Monkey V p < 0.0001). Individual Wilcoxon signed-rank tests per neuron revealed that 28/123 neurons for Monkey U, and 47/145 neurons for Monkey U showed significant changes in activity at this epoch.}
\label{fig:juice_delivery}
\end{figure}

\begin{figure}[p] 
\centering    
\includegraphics[width=1.0\textwidth]{Chapter4/Figs/redone_figs/juice_win_fractal_fig}
\caption[Averaged responses over all trials for putative dopaminergic neurons which showed significant changes in activity upon payment of the auctioned good.]{\textbf{Averaged responses over all trials for putative dopaminergic neurons which showed significant changes in activity upon payment of the auctioned good.} Neurons were selected by individual Wilcoxon signed-rank test for the difference in activity between win and lose trials. The activity of these neurons across every recorded trial is shown for trials in which the monkey won the auction (top) and lost the auction (bottom).}
\label{fig:juice_delivery_raster}
\end{figure}

When split by the fractal won on a trial (Figure \ref{fig:juice_delivery_fractal}A), both monkeys show significant increases in putative dopaminergic activity upon the payout of both the medium (Monkey U p < 0.005, Monkey V p < 0.0001) and highest (Monkey U p < 0.0001, Monkey V p < 0.0001) magnitudes of juice, though neither showed an increase in population activity when the lowest magnitude of juice was delivered (Monkey U p = 0.159, Monkey V p = 0.170; all Wilcoxon signed-rank tests). Although neither test reached significance, both monkeys showed a positive mean difference for the lowest magnitude condition (Monkey U +0.285, Monkey V +0.213 spikes/s), accounting for the apparent upward trend visible in Figure \ref{fig:juice_delivery_fractal}A). A Kruskal-Wallis test showed significant differences between the response to each reward magnitude level (Monkey U p < 0.05, Monkey V p < 0.0001) with all 3 levels showing significant rank ordering under a post hoc Dunn test for Monkey V (p < 0.0005). For Monkey U, only the difference between the high and lowest reward magnitude was significant (p < 0.0001; high-medium p = 0.144, medium-low p = 0.232).

\begin{figure}[p] 
\centering    
\includesvg[width=1.0\textwidth,inkscapelatex=false]{Chapter4/Figs/redone_figs/juice_fractal_fig}
\caption[Responses of putative dopaminergic neurons to payment of juice is ranked by the magnitude of juice received.]{\textbf{Responses of putative dopaminergic neurons to payment of juice is ranked by the magnitude of juice received.} In A) the responses of putative dopaminergic neurons to the payment of reward juice in the win condition is split by the magnitude of juice offered. The populations recorded from both monkeys show significant rank ordering of increases in activity under a Dunn test (Monkey U, high - low difference p < 0.0001; Monkey V all comparisons p < 0.0005). In B) the same split of neuronal activity is shown for the same epoch when the monkey did not receive any juice as they did not bid higher than the randomly drawn computer bid. A Kruskal-Wallis test for difference in responses to the different reward levels forfeited by the monkey (Monkey U p = 0.332; Monkey V p = 0.052).}
\label{fig:juice_delivery_fractal}
\end{figure}

In the losing condition, the inverse population activity was observed. The removal of the lowest magnitude fractal stimulus without payment of any juice led to an increase in putative dopaminergic neuron activity (Monkey U p < 0.01, Monkey V p < 0.001) but not under the removal of the medium (Monkey U p = 0.060, Monkey V p = 0.659) or highest magnitude stimulus (Monkey U p = 0.881, Monkey V p = 0.298). Neither was there a significant difference between responses to the different reward levels under a Kruskal Wallis test (Monkey U p = 0.332, Monkey V p = 0.052) though a post hoc Dunn test did show a significant difference between the medium and lowest reward magnitude activity of putative dopaminergic neurons for Monkey V (p < 0.005).

If the monkey is pleasantly surprised, or perhaps relieved, to receive their reward for completing a trial and winning an auction, the dopaminergic responses should be correlated with the subjective value the monkey places upon the potential reward in that trial. Splitting the responses to this epoch by the third of the bidspace in which the monkey had placed their bid there is a clear difference between normalised firing rates in the win (A) and lose (B) conditions in Figure \ref{fig:juice_delivery_bid}.

When the monkey loses the trial, none of the tercile splits showed a significant change in the activity of recorded dopaminergic cells (Monkey U low p = 0.358, medium p = 0.0536, high p = 0.451; Monkey V low p = 0.065, medium p = 0.069, high p = 0.731; Wilcoxon signed-rank tests). In contrast, only the activity in dopaminergic cells when Monkey V placed their bid in the lowest tercile was not significant under the same test (Monkey U low p < 0.05, medium p < 0.005, high p < 0.0001; Monkey V low p = 0.070, medium p < 0.0001, high p < 0.0001) for trials when juice was awarded for winning the auction. A Kruskal-Wallis test for differences between each group was significant for both monkeys (p < 0.0001) with only the high-medium comparison group for Monkey U not reaching significance under a post hoc Dunn test (p = 0.279; otherwise p < 0.0001). This is consistent with the visual appearance of Figure \ref{fig:juice_delivery_bid}A), in which the medium bid tercile appears to exceed the high tercile for Monkey U within the analysis window, despite the non-significant pairwise comparison.

\begin{figure}[p] 
\centering    
\includesvg[width=0.94\textwidth,inkscapelatex=false]{Chapter4/Figs/redone_figs/juice_bid_fig}
\caption[Putative dopaminergic population activity change is greatest on trials where the monkey's bid is highest.]{\textbf{Putative dopaminergic population activity change is greatest on trials where the monkey's bid is highest.} Average normalised activity of the dopaminergic neuron population (Monkey U n = 123; Monkey V n = 145 neurons) is split by tercile of the bidspace in which the monkey places their bid. A) shows responses in trials in which the monkey won the auction. The only pairwise comparison not to show significance under a Dunn test was for Monkey U comparing the change in activity on trials where the monkey bid in the highest vs. middle tercile (p = 0.279). All other comparisons were significant (p < 0.0001). Only the normalised population activity for Monkey V when the bid was placed in the lowest tercile did not show a significant increase (p = 0.070). In B) the same responses are shown for trials where the monkey did not win the auction. No grouping showed significant deviation in normalised activity from baseline.}
\label{fig:juice_delivery_bid}
\end{figure}

A simple linear regression of the form 

\begin{align}
\text{Firing Rate | Result} &\sim \beta_0 + \beta_1(\text{Monkey Bid}) + \varepsilon \label{eq:firing_rate_mbid_regression}
\end{align}

shows no significance for the loss condition for either monkey (Monkey U, p = 0.583, adjusted $R^2$ = -0.03; Monkey V, p = 0.29, adjusted $R^2$ = 0.01; Figure \ref{fig:juice_delivery_bid_regression}B). When the monkey did win the auction trial (Figure \ref{fig:juice_delivery_bid_regression}A), the simple linear model was not significant for Monkey U (p = 0.47, adjusted $R^2$ = -0.02), but it was significant for Monkey V (p < 0.005, adjust $R^2$ = 0.33). These parameters are summarised in Table \ref{tab:juice_delivery_regression_results}.

\begin{figure}[p]
\centering    
\includesvg[width=0.95\textwidth,inkscapelatex=false]{Chapter4/Figs/Vector/profit_regression_fig}
\caption[Regression of activity of putative dopaminergic cells to juice payment against monkey bid.]{\textbf{Regression of activity of putative dopaminergic cells to juice payment against monkey bid.} A regression (Equation \ref{eq:firing_rate_mbid_regression}) of the activity of putative dopaminergic cells which showed significant changes in activity upon the payment of juice (see \ref{fig:juice_delivery_raster}) is shown for both the win (top, A) and lose (bottom, B) outcomes. For neither monkey was activity at this epoch where no juice was paid significantly correlated with the bid the monkey had made for the auctioned good (Monkey U p = 0.583, adjusted $R^2$ = -0.033; Monkey V p = 0.294, adjusted $R^2$ = 0.007). In trials where the monkey's bid was greater than the randomly drawn computer bid and the monkey received a payout of juice, only Monkey V showed significant correlation (p < 0.005, adjusted $R^2$ = 0.333) between normalised neuronal activity and preceding monkey bid (Monkey U p = 0.474, adjusted $R^2$ = -0.022).}
\label{fig:juice_delivery_bid_regression}
\end{figure}

\begin{table}[htbp]
\centering
\begin{tabular}{lllccc}
\hline
\textbf{Monkey} & \textbf{Trial Outcome} & \textbf{Intercept} & \textbf{Slope} & \textbf{Adj} $R^2$ & $p$\textbf{-value} \\
\hline
Monkey U & Monkey Win & $0.28 \pm 0.08$ & $-0.09 \pm 0.13$ & -0.022 & 0.474 \\
Monkey U & Monkey Lose & $0.07 \pm 0.08$ & $-0.08 \pm 0.15$ & -0.033 & 0.583 \\
\hline
Monkey V & Monkey Win & $0.61 \pm 0.17$ & $0.88 \pm 0.27$ & 0.333 & $3.60 \times 10^{-3}$ \\
Monkey V & Monkey Lose & $0.28 \pm 0.08$ & $-0.15 \pm 0.14$ & 0.007 & 0.294 \\
\hline
\end{tabular}
\caption[Regression coefficients for the relationship between preceding monkey bid and normalised cell activity at juice payout epoch.]{\textbf{Regression coefficients for the relationship between preceding monkey bid and normalised cell activity at juice payout epoch.} For all regression captured in Equation \ref{eq:firing_rate_mbid_regression} coefficients are shown. The only significant correlation was seen when Monkey V received their juice payment having won an auction trial (see \ref{fig:juice_delivery_bid_regression}).}
\label{tab:juice_delivery_regression_results}
\end{table}

We also fitted a mixed-effects equivalent of Equation \ref{eq:firing_rate_mbid_regression}, following the same structure as Equation \ref{eq:lme_general}, separately for the win and lose conditions. This showed no significant relationship between normalised firing rate and monkey bid for either monkey when losing (Monkey U: $\beta = -0.0276 \pm 0.175$, $p = 0.875$; Monkey V: $\beta = -0.147 \pm 0.116$, $p = 0.205$), consistent with the naive regression above. When winning, however, the mixed-effects model diverges from the naive result: Monkey U remained non-significant ($\beta = -0.0855 \pm 0.144$, $p = 0.554$), and Monkey V, the only significant result in the naive analysis (Table \ref{tab:juice_delivery_regression_results}), was no longer significant once cell-level clustering was accounted for ($\beta = 0.307 \pm 0.188$, $p = 0.103$).

Notably, this null result sits in contrast to the win/lose reveal analysis above (Table \ref{tab:lme_wl_regression_results}). In both cases we would expect no correlation of dopamine activity with the predictor variable as neither the juice magnitude given a winning trial, nor the water payout given a losing trial is related to a subject's bid, and thus the signal of their utility for either good. Correcting for cell-level clustering had opposite effects in the two cases, revealing a significant relationship where none was previously found for the computer bid on loss trials, but removed one that appeared under the naive regression here. Rather than uniformly inflating or deflating significance the methodological change reveals a theoretically unexpected effect in one case and removes a likely false positive in the other.

\begin{table}[htbp]
\centering
\begin{tabular}{lllcc}
\hline
\textbf{Monkey} & \textbf{Trial Outcome} & \textbf{Intercept} & \textbf{Slope} & $p$\textbf{-value} \\
\hline
Monkey U & Monkey Win & $0.274 \pm 0.113$ & $-0.0855 \pm 0.144$ & $5.54 \times 10^{-1}$ \\
Monkey U & Monkey Lose & $0.00919 \pm 0.0904$ & $-0.0276 \pm 0.175$ & $8.75 \times 10^{-1}$ \\
\hline
Monkey V & Monkey Win & $1.05 \pm 0.301$ & $0.307 \pm 0.188$ & $1.03 \times 10^{-1}$ \\
Monkey V & Monkey Lose & $0.281 \pm 0.0968$ & $-0.147 \pm 0.116$ & $2.05 \times 10^{-1}$ \\
\hline
\end{tabular}
\caption[Mixed-effects model estimates for the relationship between preceding monkey bid and normalised cell activity at juice payout epoch.]{\textbf{Mixed-effects model estimates for the relationship between preceding monkey bid and normalised cell activity at juice payout epoch.} Results from the linear mixed-effects model described in Equation \ref{eq:lme_general}, fitted separately per monkey and trial outcome, with a random intercept per cell. Unlike the naive regression (Table \ref{tab:juice_delivery_regression_results}), no condition shows a significant relationship.}
\label{tab:lme_juice_delivery_regression_results}
\end{table}

\subsection{Dopaminergic cell activity is correlated to the magnitude of budget remaining}

At the end of each trial, whatever water budget has not been spent by the animal is paid out. Given an animal has a maximum possible bid and pays a maximum second price which is less than this, the monkey is guaranteed at least \textit{some} water in this epoch, which should be predicted at the moment of the computer bid reveal. The amount of water to be 'spent' by the monkey if they win the trial is shown on the task screen by a reversal of the gradient of the bidding space below the losing computer bid; upon payment of the budget, both this entire bidding space, the computer bid, and the monkey's bid visual stimuli disappear from the task screen leaving it empty in preparation for the next trial.

As with the receipt of the juice, even though the payout of the water can be entirely predicted at the moment that the computer's bid is revealed, the putative dopaminergic neurons we recorded showed a marked increase in activity both when the monkey loses the auction trial (and received the full budget compliment 1.2ml of water) \textit{and} when they won the trial (and received 1.2ml of water minus the proportion spent paying out the computer's 'second price.'). The normalised firing rate of these neurons relative to the delivery of the water is shown in Figure \ref{fig:water_delivery_fig}.

\begin{figure}[htbp!] 
\centering    
\includesvg[width=1.0\textwidth,inkscapelatex=false]{Chapter4/Figs/redone_figs/water_payout_fig}
\caption[Putative dopaminergic neurons show increases in firing rate when the leftover water budget is delivered.]{\textbf{Putative dopaminergic neurons show increases in firing rate when the leftover water budget is delivered.} In both the win (purple) and lose (black) condition, the putative dopaminergic neurons for both monkeys (Monkey U, left, n = 123 neurons; Monkey V, right, n = 145 neurons) show increases in firing rate at the final task epoch. In both these conditions, at least some water is delivered to the monkey. When the monkey loses the trial, the monkey receives the full 1.2ml of budget water. When the monkey bids higher than the randomly drawn computer bid and wins the trial, they receive 1.2ml-the proportion of the budget bid by the computer. This can be anything from a minuscule amount (if the monkey bid the full budget and the computer only very slightly less) to almost the full 1.2ml.}
\label{fig:water_delivery_fig}
\end{figure}

In comparison to the previous two (reveal of the trial outcome and payment of any juice reward) epochs, the greatest increase in firing rate is seen when the monkey loses the trial, suggesting that the change in firing rate is proportional to the magnitude of budget liquid received. To investigate, we performed binned linear regressions of the change in firing rate against both the proportion of total budget received ('budget water remaining') and the proportion of the total budget that was 'profit' (defined as the monkeys bid minus the randomly drawn computer bid).

In Figure \ref{fig:water_delivery_regression_fig}, only the total magnitude of the water delivered is significant for both monkeys (top; Monkey U p = 0.022, Monkey V p < 0.0001) in linear regressions of the form:

\begin{align}
\text{Firing Rate | Result} &\sim \beta_0 + \beta_1(\text{Budget Remaining}) + \varepsilon \label{eq:firing_rate_budget_regression1}
\end{align}
\begin{align}
\text{Firing Rate | Result} &\sim \beta_0 + \beta_1(\text{Budget Profit}) + \varepsilon \label{eq:firing_rate_budget_regression2}
\end{align}


The amount of budget that the monkey gained in profit per trial was significant for neither animal (Monkey U; p = 0.295; Monkey V p = 0.586. We included only trials in which the monkey won the auction to remove any effects of a large, but disappointing, budget payout in a loss trial. These results are summarised in Table \ref{tab:water_delivery_regression_results}.

\begin{figure}[htbp!] 
\centering    
\includesvg[width=1.0\textwidth,inkscapelatex=false]{Chapter4/Figs/redone_figs/water_regression_fig}
\caption[Putative dopaminergic activity is correlated with the amount of water returned to the monkey.]{\textbf{Putative dopaminergic activity is correlated with the amount of water returned to the monkey.} For both Monkey V and Monkey U, there was significant linear correlation (see \ref{tab:water_delivery_regression_results}) between the change in firing rate of putative dopaminergic neurons and the total amount of budget water returned to the monkey. Only trials where the monkey won were included. In contrast, no significant relationship was seen between the change in firing rate and the profit (the monkey's bid - the computer's bid) that the water delivered represented on trials where the monkey won the auction.}
\label{fig:water_delivery_regression_fig}
\end{figure}

\begin{table}[htbp]
\centering
\begin{tabular}{lllccc}
\hline
\textbf{Monkey} & \textbf{Ind. Variable} & \textbf{Intercept} & \textbf{Slope} & \textbf{Adj} $R^2$ & $p$\textbf{-value} \\
\hline
Monkey U & Budget Remaining & $0.04 \pm 0.03$ & $0.12 \pm 0.05$ & 0.173 & 0.022 \\
Monkey U & Budget Profit & $0.10 \pm 0.04$ & $0.06 \pm 0.06$ & 0.006 & 0.295 \\
\hline
Monkey V & Budget Remaining & $0.14 \pm 0.03$ & $0.37 \pm 0.06$ & 0.614 & $2.18 \times 10^{-6}$ \\
Monkey V & Budget Profit & $0.38 \pm 0.06$ & $-0.06 \pm 0.11$ & -0.031 & 0.586 \\
\hline
\end{tabular}
\caption[Regression coefficients for the relationship between budget payout and normalised cell activity.]{\textbf{Regression coefficients for the relationship between budget payout and normalised cell activity.} Full statistics for the linear regressions described in Equations \ref{eq:firing_rate_budget_regression1} and \ref{eq:firing_rate_budget_regression2} are shown. Significant relationships were only found in the regression of dopaminergic activity against the magnitude of budget remaining and delivered to the each monkey.}
\label{tab:water_delivery_regression_results}
\end{table}

We also fitted mixed-effects equivalents of Equations \ref{eq:firing_rate_budget_regression1} and \ref{eq:firing_rate_budget_regression2}, following the same structure as Equation \ref{eq:lme_general}. This confirmed a significant positive relationship between normalised firing rate and budget remaining for both monkeys ($\beta = 0.122 \pm 0.0598$, $p = 0.042$ for Monkey U; $\beta = 0.305 \pm 0.0522$, $p = 5.15 \times 10^{-9}$ for Monkey V), and no significant relationship with budget profit for either monkey ($\beta = 0.0125 \pm 0.0601$, $p = 0.835$ for Monkey U; $\beta = 0.00376 \pm 0.0522$, $p = 0.943$ for Monkey V), consistent with the naive regression above in both direction and significance.

\begin{table}[htbp]
\centering
\begin{tabular}{lllcc}
\hline
\textbf{Monkey} & \textbf{Ind. Variable} & \textbf{Intercept} & \textbf{Slope} & $p$\textbf{-value} \\
\hline
Monkey U & Budget Remaining & $0.0278 \pm 0.0512$ & $0.122 \pm 0.0598$ & $4.22 \times 10^{-2}$ \\
Monkey U & Budget Profit & $0.105 \pm 0.0376$ & $0.0125 \pm 0.0601$ & $8.35 \times 10^{-1}$ \\
\hline
Monkey V & Budget Remaining & $0.153 \pm 0.101$ & $0.305 \pm 0.0522$ & $5.15 \times 10^{-9}$ \\
Monkey V & Budget Profit & $0.341 \pm 0.0979$ & $0.00376 \pm 0.0522$ & $9.43 \times 10^{-1}$ \\
\hline
\end{tabular}
\caption[Mixed-effects model estimates for the relationship between budget payout and normalised cell activity.]{\textbf{Mixed-effects model estimates for the relationship between budget payout and normalised cell activity.} Results from the linear mixed-effects model described in Equation \ref{eq:lme_general}, fitted separately per monkey and predictor, with a random intercept per cell. Consistent with the naive regression (Table \ref{tab:water_delivery_regression_results}), a significant relationship is found only for budget remaining for both monkeys, not the budget profit.}
\label{tab:lme_water_delivery_regression_results}
\end{table}

We mentioned when discussing the payout juice that it is not surprising that the putative dopaminergic neurons recorded in this project show these stereotyped reward prediction errors despite the animal having perfect knowledge of what they will receive. For the RPEs shown in the water delivery epoch, we can think of at least two reasons why these would persist:

Firstly, despite the animal in theory having perfect knowledge of the budget they will receive upon the reveal of the trial outcome, there is a long gap of 3 seconds until the water is actually released. In addition, halfway through this time any juice won is given to the monkey which may break any concentration about the water budget to be paid out.

Secondly, on any trial it is very unlikely (unless the monkey loses sequentially) that the budget leftover is the same. Even though the amount left is signalled by the horizontal computer bid bar and reversal of the shaded budget pattern, the monkey has to judge the amount that this reversal corresponds to by eye, which is unlikely to be perfectly accurate. Therefore, there will always be \textit{some} error on the monkey's prediction of water received. On top of this error, the liquid flows to the monkey via gravitational force after opening a solenoid valve, which will always result in small fluctuations of liquid for the same amount of time the tap is left open.

A one-sided Wilcoxon signed-rank test found that 43/123 putative dopaminergic neurons recorded from Monkey U, and 66/145 putative dopaminergic neurons recorded from Monkey V showed a significant increase in activity following the delivery of the water budget. This was tested only on trials in which the monkey lost the auction to remove any effect of the size of the budget returned.

A full table of the number of putative dopaminergic neurons recorded per monkey is shown in Table \ref{tab:all_cells_significance}

\begin{table}[htbp]
\centering
\label{tab:monkey_data}
\begin{tabular}{lrrr}
\hline
\textbf{Description} & \textbf{Monkey U} & \textbf{Monkey V} & \textbf{Total} \\
\hline
All cells & 236 & 259 & 495 \\
Putative dopaminergic & 123 & 145 & 268 \\
Other & 113 & 114 & 227 \\
\hline
\textbf{Responsiveness} & & & \\
Free reward & 80 & 33 & 113 \\
Fixation cross & 82 & 62 & 144 \\
Fractal offer & 31 (of 102) & 41 (of 121) & 72 (of 223) \\
Single fractal offer & 3 (of 21) & 7 (of 24) & 10 (of 45) \\
Monkey bid & 32 & 41 & 73 \\
Win/lose & 72 & 46 & 118 \\
Juice delivery & 28 & 47 & 75 \\
Water delivery & 43 & 66 & 75 \\
\hline
\end{tabular}
\caption[Summary of all cells recorded from Monkey V and Monkey U.]{\textbf{Summary of all cells recorded from Monkey V and Monkey U.} Table shows counts of all cells recorded in this project. In total 123 (Monkey U) and 145 (Monkey V) putative dopaminergic cells were recorded. The number of cells which showed significant changes in activity to the delivery of unpredicted water ('Free reward'), the onset of the trial ('Fixation cross'), the reward offer magnitude ('Fractal offer'), the forthcoming bid made by the monkey ('Monkey bid'), the outcome of the trial ('Win/lose') and the delivery of juice reward or water budget are shown. All denominators for responsive cells are the total number of putative dopaminergic neurons except for the fractal offer where the number of recorded cells under each paradigm is shown.}
\label{tab:all_cells_significance}
\end{table}